% Modernised reading — fonds Alexandre Grothendieck, shelfmark 19, pages 1-95.
%
% This is an INTERPRETATION, not a transcription. It re-reads the pages in
% today's notation and vocabulary, states the hypotheses the manuscript leaves
% implicit, and drops the critical apparatus so that the mathematics reads
% straight through. Everything the brackets and underlines carried — a doubtful
% reading, a notation he used differently, a missing passage — moves into a
% footnote.
%
% It is held to being mathematically correct as it stands, not merely faithful
% to the page. Where the manuscript is loose or mistaken, the true statement is
% what appears here, and the footnote says what the page has. In any dispute of
% substance the transcription governs: whoever cites Grothendieck must cite the
% batch-NN.fr.tex files.
%
% Scope: the folder was read WHOLE before any of this was written — all five
% batches, pages 1-95, in one pass — and it is written as ONE file for the
% whole shelfmark, which is this project's layout: `19.modern.tex` covers the
% five batches of transcription, exactly as 108.modern.tex, 115.modern.tex and
% 161-3.modern.tex cover theirs.
%
% One notation is fixed here for the whole folder and held to throughout, in
% the "Conventions" section: unit, counit and comultiplication get modern
% letters, because the manuscript's own letters are NOT consistent from one
% bundle to the next (alpha and beta are swapped between pages 3 and 70).
%
% Four corrections are substantive rather than cosmetic, and each is footnoted
% where it occurs:
%   - h_B is covariant in B, not contravariant as page 14 says;
%   - the equivalence of page 15 needs the generator U to be small projective;
%     with a bare generator one gets Gabriel-Popescu, not an equivalence;
%   - page 60 gives i* a kernel of i_! where the symmetry requires j_*;
%   - page 70 says "g adjoint à gauche de f" where its own arrows say otherwise.
%
% One convention is chosen openly rather than silently, at the head of the
% recollement section: Grothendieck's i is the OPEN immersion and his j the
% CLOSED one, the exact reverse of today's usage. His letters are kept, so that
% this reading can be laid beside the transcription formula by formula.
%
% Pass: Opus 5 (claude-opus-5), 2026-08-22 — first pass, unchecked against the
% pages by a human.
\documentclass[11pt,a4paper]{article}
% Le préambule commun à toutes les transcriptions du fonds Grothendieck.
%
% Il tient en une idée : **l'incertitude se compose, elle ne se résout pas.**
% Une transcription qui lisse un mot douteux a détruit la seule chose pour
% laquelle on la faisait — un lecteur doit pouvoir savoir, à chaque endroit,
% ce qui était sur le feuillet et ce qui a été ajouté. Les six macros qui
% suivent sont donc l'appareil critique, et non de la décoration.
%
% Les mêmes commandes sont reconnues par `scripts/render.mjs`, qui produit la
% vue de lecture du site. Une macro ajoutée ici doit l'être aussi là-bas,
% faute de quoi le rendu échouera bruyamment — ce qui est voulu : un rendu
% silencieusement faux est pire qu'un rendu absent.

\ProvidesPackage{grothendieck}[2026/08/08 Transcriptions du fonds Grothendieck]

\usepackage{amsmath,amssymb,amsthm}
\usepackage{mathrsfs}
\usepackage{stmaryrd}
% Les diagrammes commutatifs sont partout dans ce fonds : c'est la langue
% même de la géométrie algébrique de Grothendieck, et une transcription qui
% les rendrait en prose ne transcrirait rien.
\usepackage{tikz-cd}
% Français par défaut — c'est la langue des feuillets — mais l'anglais est
% chargé aussi : la traduction et le résumé sont des documents anglais, et la
% typographie française (espaces avant les deux-points) y serait une faute.
% Ces documents basculent par \selectlanguage{english} en tête de corps.
\usepackage[english,french]{babel}
\usepackage{fontspec}
\usepackage{geometry}
\geometry{a4paper,margin=25mm,marginparwidth=28mm}
\usepackage{marginnote}
\usepackage{xcolor}
% ulem, not soul. `soul`'s \st and \ul are built on a character-by-character
% scan that XeLaTeX's font machinery breaks: the document fails with an
% undefined control sequence rather than degrading. ulem does the same two jobs
% and is Unicode-engine safe. `normalem` stops it hijacking \emph, which the
% transcriptions use for Grothendieck's own emphasis.
\usepackage[normalem]{ulem}
\usepackage{hyperref}
\hypersetup{colorlinks=true,linkcolor=black,urlcolor=black,pdfborder={0 0 0}}

% --- Métadonnées du lot -----------------------------------------------------
% Reprises telles quelles par le script de rendu, qui les affiche en tête de
% la vue de lecture. Elles fixent ce que le fichier prétend couvrir : sans
% elles, une transcription flotte sans point d'attache dans un fonds de seize
% mille feuillets.

\newcommand{\folder}[1]{\gdef\grothFolder{#1}}
\newcommand{\batch}[1]{\gdef\grothBatch{#1}}
\newcommand{\leaves}[2]{\gdef\grothFirst{#1}\gdef\grothLast{#2}}
\newcommand{\foldertitle}[1]{\gdef\grothFoldertitle{#1}}
\gdef\grothFolder{?}\gdef\grothBatch{?}\gdef\grothFirst{?}\gdef\grothLast{?}\gdef\grothFoldertitle{}

% --- Le résumé --------------------------------------------------------------
%
% Il ouvre la lecture modernisée, sous le titre « Résumé », et il est composé
% en petit. Ce n'est pas un ornement : c'est le seul passage du document qui
% ne soit pas des mathématiques, et un lecteur doit pouvoir le reconnaître
% avant de l'avoir lu — puis le sauter s'il connaît déjà le sujet, ou s'y
% arrêter s'il ne le connaît pas. Le filet qui le suit marque la reprise du
% corps.

\newenvironment{resume}
  {\small\setlength{\parskip}{0.45em}}
  {\par\medskip\hrule\medskip}

% --- Mots-clés ---------------------------------------------------------------
%
% En anglais, en clôture du résumé de la lecture modernisée : le vocabulaire
% moderne sous lequel chercher ce que les feuillets construisent. Le manifeste
% du site (`npm run manifest`) les extrait pour étiqueter la cote entière —
% cette ligne est la source unique des tags, il n'y a pas de fichier à côté.
\newcommand{\keywords}[1]{\par\smallskip\noindent
  {\footnotesize\textsc{Keywords} --- \textit{#1}}\par}

% --- L'appareil critique ----------------------------------------------------

% Le numéro de feuillet, en marge. C'est l'ancre qui permet de régler un
% désaccord sur une formule en regardant *un* feuillet plutôt qu'une chemise
% de six cents. Le site s'en sert aussi pour tourner les pages du fac-similé
% quand on fait défiler la transcription.
\newcommand{\leaf}[1]{%
  \marginnote{\footnotesize\color{gray!70}\textsf{#1}}%
  \label{leaf:#1}\ignorespaces}

% Illisible : on ne devine pas. Un mot inventé qui se lit comme les autres est
% le pire résultat possible.
\newcommand{\ill}{\textcolor{red!60!black}{[\,\ldots\,]}}

% Lecture douteuse : le mot est donné, et signalé comme incertain.
\newcommand{\uncertain}[1]{\uline{#1}}

% Ajout de l'éditeur — une lettre manquante, un mot sous-entendu.
\newcommand{\add}[1]{\textcolor{blue!50!black}{[#1]}}

% Biffé par Grothendieck lui-même. Ce qu'il a rayé fait partie du document :
% une rature dit souvent où la pensée a changé de direction.
\newcommand{\struck}[1]{\sout{#1}}

% Note du transcripteur, sur l'état matériel du feuillet.
\newcommand{\note}[1]{\par\smallskip{\small\color{gray!60!black}\textit{[#1]}}\par\smallskip}

% Note marginale de Grothendieck, distincte du corps du texte.
\newcommand{\marginal}[1]{\par\smallskip\noindent\hspace{1em}%
  {\small\color{gray!60!black}#1}\par\smallskip}

% --- Titre ------------------------------------------------------------------

\AtBeginDocument{%
  \begin{center}
    {\large\bfseries Fonds Alexandre Grothendieck --- cote n\textsuperscript{o}~\grothFolder}\\[2pt]
    {\itshape\grothFoldertitle}\\[4pt]
    {\small feuillets \grothFirst--\grothLast\ (lot \grothBatch)}\\[2pt]
    {\footnotesize\color{gray!60!black}Transcription établie d'après le fac-similé numérisé
     par l'Université de Montpellier.\\
     Les lectures incertaines sont soulignées, les illisibles notées [\,\ldots\,],
     les ajouts entre crochets.}
  \end{center}
  \vspace{1em}}

\endinput


\folder{19}
\pages{1}{95}
\dating{[à partir de 1958-1973]}
\watermark{Édition de démonstration\\interprétation personnelle de l'œuvre}
\foldertitle{Topos : notes manuscrites (s.d.) — lecture modernisée du dossier entier}

\begin{document}

\section*{Résumé}

\begin{resume}
La chemise porte deux mots de l'inventaire : « Topos, notes manuscrites ».
Dedans, neuf liasses qui n'ont pas l'air d'aller ensemble — des coalgèbres, de
la cohomologie relative, des notes de rédaction pour un exposé de SGA 4, des
catégories de fractions, un théorème de recollement, une liste de questions
sans réponses, des catégories fibrées. Ce n'est pas un texte : c'est un dossier
de travail. Mais une seule question le traverse d'un bout à l'autre, et il vaut
la peine de l'avoir en tête avant d'ouvrir n'importe laquelle de ses pages.

Voici cette question. Presque toutes les constructions des mathématiques
\emph{oublient} quelque chose. D'un groupe on ne retient que l'ensemble
sous-jacent ; d'un faisceau sur un espace, sa restriction à un ouvert ; d'un
module sur un anneau, le même module vu sur un sous-anneau. Chaque fois, on a
un procédé $f$ qui va d'une collection d'objets riches vers une collection
d'objets plus pauvres, et chaque fois on perd de l'information. Combien
exactement ? Et surtout : ce qui a été perdu peut-il être \emph{rangé quelque
part}, sous une forme telle qu'à partir de l'objet pauvre et de ce rangement on
reconstitue l'objet riche, exactement et sans reste ?

L'image juste est celle du recollement. Pour connaître une fonction sur un
espace, il suffit de la connaître sur chaque morceau d'un recouvrement, à
condition de savoir en plus que les valeurs coïncident sur les recouvrements
deux à deux : ces conditions de coïncidence sont la « donnée de recollement »,
et sans elles les morceaux ne se recollent pas. Le pari du dossier est que
cette situation-là est universelle. Dès que le procédé $f$ admet un
\emph{adjoint} $g$ — un procédé qui, à un objet pauvre, associe le meilleur
objet riche qui s'envoie dessus —, le composé $\varphi = fg$ mesure, sur les
objets pauvres eux-mêmes, ce que la remontée puis la redescente font subir. Et
une « donnée de recollement » sur un objet $X$ n'est rien d'autre qu'une flèche
$X \to \varphi(X)$ obéissant à deux règles de compatibilité. Grothendieck
appelle ces objets-avec-flèche des \emph{tapis cartésiens} ; on dit aujourd'hui
des \emph{coalgèbres}.

Le théorème qui ouvre le dossier dit quand le recollement est effectif : quand
la catégorie des coalgèbres est, à équivalence près, la catégorie riche de
départ. Deux conditions suffisent, et elles sont nécessaires — que $f$ ne
déclare isomorphes que des objets qui le sont déjà, et qu'il respecte assez de
noyaux. C'est le théorème de comonadicité, celui qu'on appelle aujourd'hui
théorème de Beck ; la page en porte le nom.

Tout le reste met ce mécanisme à l'épreuve, sur des objets de plus en plus
gros, et c'est ce qui donne au dossier son unité.

Il sert d'abord à fabriquer des résolutions — la \emph{construction standard}
de Godement, où l'on itère $\varphi$ pour obtenir un objet cosimplicial, dont
tout l'intérêt est d'être trivial du côté où il a été construit. Il sert
ensuite à faire de l'algèbre homologique \emph{relative} : on décrète exactes
non pas toutes les suites exactes, mais celles que $f$ scinde, et toute la
théorie des foncteurs dérivés se refait avec cette notion-là. Les objets
injectifs relatifs se reconnaissent alors sans peine — ce sont exactement les
facteurs directs des objets $T(N)$ — et le foncteur $C^{0}$ de Godement, celui
qui à un faisceau associe le produit de ses fibres, apparaît comme un cas
particulier de la construction générale.

Il sert à \emph{localiser}, c'est-à-dire à rendre inversible une famille de
flèches, et à reconnaître les objets que la localisation ne bouge pas : c'est
la théorie de Gabriel, et le dossier en donne les conditions d'existence du
foncteur section. Il sert à \emph{recoller} un ouvert et son fermé
complémentaire — le théorème que la chemise attribue à Michael Artin et Pierre
Cartier, et dont sort le tableau de six foncteurs qu'on appelle aujourd'hui un
recollement. Il sert à dire ce qu'est un \emph{foncteur d'oubli de structure} :
un foncteur fidèle et transportable, ce qui se lit entièrement sur ses fibres,
et l'on retrouve au passage que se donner une catégorie au-dessus de $B$ à
fibres discrètes revient à se donner un préfaisceau sur $B$. Il sert enfin à
décrire une \emph{famille} de catégories paramétrée par une autre — les
catégories fibrées, sur lesquelles le dossier s'achève, et dont le topos total
se calcule comme la catégorie des sections cartésiennes, c'est-à-dire des
familles recollées.

Deux liasses font exception et sont d'une autre nature : des notes de rédaction
pour les exposés IX et XII de SGA 4, sur les faisceaux constructibles et la
pureté. Elles sont écrites vite, très petit, et une grande part de leur prose
de liaison est illisible ; on n'a pas cherché à la rendre lisible, et les
silences sont dits comme tels. Une chose y vaut d'être vue malgré tout : le
geste qui y revient — ramener une propriété de tous les objets à une propriété
d'une famille génératrice, puis dévisser — est le réflexe central du dossier et
ne dépend pas du sujet.

On ne trouvera nulle part ici de rédaction. On trouve un homme qui pose ses
questions par écrit, qui numérote, qui raye, qui écrit « faux en général ! »
dans la marge en face de sa propre proposition, et qui laisse en fin de dossier
onze problèmes ouverts avec, en regard de deux d'entre eux, un « O.K. » et un
« non » notés plus tard, d'une autre encre. C'est un état de travail, et il est
donné tel quel.

\keywords{comonad, Beck comonadicity theorem, coalgebra, standard resolution, Gabriel-Popescu theorem, relative homological algebra, allowable class, effaceable functor, constructible sheaf, cohomological dévissage, Chow's lemma, Gabriel localization, quotient category, localizing subcategory, recollement, Artin gluing, torsion theory, six-functor formalism, extra degeneracy, discrete fibration, category of elements, adjoint string, transport of structure, fibred site, total topos, cartesian section, pseudofunctor, opfibration, 2-limit}
\end{resume}
\pagerange{3}{3}
\section*{Conventions, valables pour tout le dossier}

Les lettres changent d'une liasse à l'autre : ce que la page 3 note $\beta$, la
page 70 le note $\alpha$, et inversement. On fixe donc ici une notation, et on
s'y tient dans les cinq lectures modernisées de ce dossier.\footnote{Le
dictionnaire est le suivant. Page 3 : $\beta$ est l'unité, $\alpha$ la coünité,
$\lambda$ la comultiplication. Page 13 : $\phi$ est l'unité de la monade
$E = GF$ et $\lambda$ sa multiplication ; $\phi'$ et $\lambda'$ sont la coünité
et la comultiplication de la comonade $E' = FG$. Page 70 : $\alpha$ est
l'unité, $\beta$ la coünité, $\lambda$ la comultiplication — soit l'échange de
$\alpha$ et $\beta$ par rapport à la page 3. C'est parce que les lettres du
manuscrit ne sont pas cohérentes d'une liasse à l'autre qu'on ne les a pas
reprises.}

Une \emph{adjonction} $f \dashv g$ est la donnée de deux foncteurs
$f : A \to B$ et $g : B \to A$ et d'une bijection
\[
  \mathrm{Hom}_{B}\bigl( f(Y), X \bigr) \;\simeq\;
  \mathrm{Hom}_{A}\bigl( Y, g(X) \bigr) ,
\]
naturelle en $Y$ et $X$ ; on dit que $f$ est adjoint à gauche et $g$ adjoint à
droite. Elle équivaut à la donnée de deux transformations naturelles, l'unité
$\eta : \mathrm{id}_{A} \to gf$ et la coünité
$\varepsilon : fg \to \mathrm{id}_{B}$, soumises aux deux identités
triangulaires.

Le composé $\varphi = fg : B \to B$ est alors une \emph{comonade} : il vient
avec la coünité $\varepsilon : \varphi \to \mathrm{id}_{B}$ et une
comultiplication
\[
  \delta \;=\; f * \eta * g \;:\; \varphi \longrightarrow \varphi^{2} ,
\]
et ces deux flèches satisfont deux axiomes, un axiome de coünité et un axiome
de coassociativité, exprimés par la commutation de

\begin{tikzcd}[column sep=large]
\varphi \arrow[r, "\delta"] \arrow[rr, bend right=32, "\mathrm{id}_{\varphi}"'] &
\varphi^{2} \arrow[r, bend left=18, "\varphi * \varepsilon"] \arrow[r, bend right=18, "\varepsilon * \varphi"'] &
\varphi
\end{tikzcd}

\begin{tikzcd}[column sep=large]
\varphi \arrow[r, "\delta"] &
\varphi^{2} \arrow[r, bend left=18, "\varphi * \delta"] \arrow[r, bend right=18, "\delta * \varphi"'] &
\varphi^{3}
\end{tikzcd}

Ce sont les \emph{conditions de Godement}, du nom sous lequel le manuscrit les
cite.\footnote{Elles apparaissent dans la \emph{Topologie algébrique et théorie
des faisceaux} de Godement, 1958, attachées à la « construction standard ». Le
mot \emph{comonade} et la théorie générale sont postérieurs ; le manuscrit ne
les emploie pas.} Symétriquement, $\psi = gf : A \to A$ est une \emph{monade},
d'unité $\eta$ et de multiplication $\mu = g * \varepsilon * f$.

\pagerange{3}{4}
\section*{« Tapis cartésiens » : la comonade d'une adjonction (pages 1 à 11)}

Le premier texte du dossier occupe les pages 3 à 11 ; la chemise, page 1, lui
donne son titre et donne aussi celui du texte suivant.

\subsection*{Les coalgèbres}

Soit $\varphi$ une comonade sur $B$, de coünité $\varepsilon$ et de
comultiplication $\delta$. Une \emph{$\varphi$-coalgèbre} est un couple
$(X, u)$ formé d'un objet $X$ de $B$ et d'une flèche
\[
  u : X \longrightarrow \varphi(X)
\]
telle que les deux diagrammes suivants commutent :

\begin{tikzcd}[column sep=large]
X \arrow[r, "u"] & \varphi(X) \arrow[r, "\varepsilon(X)"] & X
\end{tikzcd}

\begin{tikzcd}[column sep=large]
X \arrow[r, "u"] &
\varphi(X) \arrow[r, bend left=18, "\varphi(u)"] \arrow[r, bend right=18, "\delta(X)"'] &
\varphi^{2}(X)
\end{tikzcd}

soit $\varepsilon(X) \circ u = \mathrm{id}_{X}$ (coünité) et
$\delta(X) \circ u = \varphi(u) \circ u$ (coassociativité). Les morphismes de
coalgèbres sont les flèches de $B$ compatibles avec les $u$. On note
$\varphi\text{-}\mathbf{Coalg}$ cette catégorie ; le manuscrit la note
$K(\varphi, \alpha, \lambda)$ et appelle ses objets des \emph{tapis
cartésiens}. C'est la catégorie de co-Eilenberg--Moore de la
comonade.\footnote{Le manuscrit nomme Eilenberg--Moore une seule fois, page 19,
et à propos de la situation duale. La définition ci-dessus ne fait intervenir
ni la coünité ni la coassociativité de $\varphi$ : le manuscrit le remarque
explicitement page 4, entre parenthèses et avec un point d'exclamation. Ces
axiomes ne servent qu'à partir du moment où l'on veut construire l'adjoint $g$.}

Lorsque $\varphi = fg$ provient d'une adjonction $f \dashv g$, un objet $Y$ de
$A$ donne canoniquement une coalgèbre, à savoir $f(Y)$ munie de
\[
  f(Y) \xrightarrow{\;f * \eta\;} fgf(Y) = \varphi\bigl( f(Y) \bigr) ,
\]
et ceci est fonctoriel. On obtient le \emph{foncteur de comparaison}
\[
  h : A \longrightarrow \varphi\text{-}\mathbf{Coalg} .
\]
Toute la question est de savoir quand $h$ est une équivalence : c'est-à-dire
quand la donnée d'un objet de $A$ équivaut à celle d'un objet de $B$ muni d'une
donnée de recollement.

\pagerange{4}{4}
\subsection*{Le théorème de comonadicité}

\textbf{Théorème.} Soit $f \dashv g$ une adjonction, $\varphi = fg$, et $h$ le
foncteur de comparaison.

\begin{enumerate}
\item[a)] Si les noyaux de doubles flèches existent dans $A$ et si $f$ y
commute, alors $h$ est essentiellement surjectif.
\item[b)] $h$ est une équivalence si et seulement si $f$ est
\emph{conservatif} — il ne rend isomorphe que ce qui l'était déjà — et si, pour
toute double flèche $(u,v)$ de $A$ telle que $\mathrm{Ker}\bigl(f(u),
f(v)\bigr)$ existe, $\mathrm{Ker}(u,v)$ existe et $f$ y commute.
\end{enumerate}

La restriction portée en marge — il suffit de considérer les doubles flèches
$(u,v)$ dont l'image par $f$ a un noyau — est la condition que le manuscrit
appelle \emph{condition C}, et c'est elle qui fait la force de l'énoncé : on ne
demande rien de $A$ en général, seulement de ce que $f$ voit déjà. C'est le
théorème de comonadicité de Beck, sous sa forme précise.\footnote{La page porte
un nom propre entre parenthèses après « Théorème 1.12 », que la transcription
lit « Beck » sans certitude. La forme monadique du théorème est celle de la
thèse de Jon Beck (1967) ; la forme comonadique, qui est celle-ci, s'en déduit
par dualité. Les doubles flèches dont $f$ voit le noyau sont, dans le
vocabulaire d'aujourd'hui, les paires $f$-scindées.}

\pagerange{4}{6}
\subsection*{La réciproque : reconstruire l'adjonction}

Partons cette fois d'une comonade $(\varphi, \varepsilon, \delta)$ sur $B$,
sans supposer qu'elle provienne de quoi que ce soit. Posons
$A = \varphi\text{-}\mathbf{Coalg}$ et soit $f : A \to B$ le foncteur d'oubli,
$(X, u) \mapsto X$. On définit
\[
  g : B \longrightarrow A , \qquad
  g(X) = \bigl( \varphi(X),\; \delta(X) \bigr) ,
\]
et c'est ici, et seulement ici, que la coassociativité sert : elle dit
exactement que $\delta(X)$ est une structure de coalgèbre sur $\varphi(X)$. On
a alors, par construction,
\[
  fg = \varphi .
\]

Que $f$ soit adjoint à gauche de $g$ se vérifie à la main. Dans un sens, la
coünité $\varepsilon$ fournit
\[
  \mathrm{Hom}_{A}\bigl( (X,u), g(Y) \bigr) \longrightarrow
  \mathrm{Hom}_{B}(X, Y) , \qquad w \longmapsto \varepsilon(Y) \circ f(w) ;
\]
dans l'autre, on construit l'unité $\eta : \mathrm{id}_{A} \to gf$ en associant
à une coalgèbre $(X,u)$ le carré

\begin{tikzcd}
X \arrow[r, "u"] \arrow[d, "u"'] & \varphi(X) \arrow[d, "\varphi(u)"] \\
\varphi(X) \arrow[r, "\delta(X)"'] & \varphi^{2}(X)
\end{tikzcd}

dont la commutation est précisément l'axiome de coassociativité de la
coalgèbre : c'est donc bien un morphisme de coalgèbres
$(X,u) \to gf(X,u) = \bigl( \varphi(X), \delta(X) \bigr)$. Les deux
applications sont inverses l'une de l'autre — l'une des deux identités
triangulaires résulte de l'axiome de coünité de $\varphi$, l'autre de l'axiome
de coünité de la coalgèbre — et l'on retrouve $(\varphi, \varepsilon, \delta)$
par le procédé du paragraphe précédent.

\textbf{Conclusion.} Pour $B$ fixée, se donner une catégorie $A$ et une
adjonction $f \dashv g$ entre $A$ et $B$, avec $f$ conservatif et satisfaisant
à la condition C, revient essentiellement à se donner une comonade
$(\varphi, \varepsilon, \delta)$ sur $B$.\footnote{Le manuscrit écrit ici « la
donnée d'un foncteur $g$ dans $B$, avec les données de Godement $\alpha$,
$\lambda$ » ; c'est $\varphi$, endofoncteur de $B$, qu'il faut lire, et non
l'adjoint $g$, qui n'a plus de sens une fois $A$ oubliée.}

\pagerange{6}{6}
\subsection*{Exactitude, et le cas des topos}

Une comonade transporte des propriétés d'exactitude de $B$ à
$\varphi\text{-}\mathbf{Coalg}$, mais pas symétriquement, et c'est un point que
la page énonce trop vite.

Le foncteur d'oubli $f$ \emph{crée} toutes les limites inductives qui existent
dans $B$ : une limite inductive de coalgèbres se calcule en oubliant les
structures, et la structure sur la limite s'en déduit sans hypothèse. Pour les
limites projectives il faut une condition, et c'est $\varphi$ qui la porte :
si $\varphi$ commute aux limites projectives d'un type donné, alors ces
limites existent dans $\varphi\text{-}\mathbf{Coalg}$ et $f$ y
commute.\footnote{La page énonce deux propriétés, a) et b), toutes deux avec
$\varprojlim$. L'une des deux doit se lire $\varinjlim$ : les limites
inductives passent sans condition, les limites projectives réclament que
$\varphi$ y commute. C'est du reste ce qu'exige le théorème sur les topos qui
suit, où l'exactitude à gauche de $\varphi$ est précisément l'hypothèse.}

\textbf{Théorème.} Supposons $B$ un topos, et $\varphi$ exact à gauche. Alors
$\varphi\text{-}\mathbf{Coalg}$ est un topos, et le foncteur d'oubli $f$ est
l'image inverse $f_{0}^{*}$ d'un morphisme de topos
$f_{0} : B \to \varphi\text{-}\mathbf{Coalg}$, morphisme qui est surjectif.
Pour que $f_{0}$ soit \emph{essentiel}, il faut et il suffit que $\varphi$
commute aux limites projectives quelconques, c'est-à-dire — $\varphi$ étant
déjà exact à gauche — aux produits infinis.\footnote{La page écrit
« accessible (i.e.\ conservatif) » là où on a écrit « surjectif ». Un morphisme
de topos est une surjection exactement quand son image inverse est fidèle, ce
qui pour un foncteur exact entre topos équivaut à la conservativité : la
parenthèse de la page est donc juste, et c'est le mot qu'elle qualifie qui
n'est pas celui d'aujourd'hui. Un morphisme est dit \emph{essentiel} lorsque
son image inverse admet elle-même un adjoint à gauche ; la page suppose
d'ailleurs explicitement que $\varphi$ a un adjoint à gauche.}

C'est l'énoncé qui justifie le titre du dossier : les comonades exactes à
gauche sur un topos sont exactement ses surjections, et la théorie des tapis
cartésiens est, pour un topos, une théorie de la descente.

\pagerange{7}{11}
\subsection*{Une base produit : la matrice des $\varphi_{ji}$}

Le reste du texte examine ce que devient tout ceci quand
$B = \prod_{i \in I} B_{i}$. Le foncteur $f$ est alors une famille
$f_{i} : A \to B_{i}$, l'adjoint $g$ une famille $g_{i} : B_{i} \to A$ avec
$g\bigl( (X_{i}) \bigr) = \prod_{i} g_{i}(X_{i})$, et la comonade $\varphi$ se
décompose. En posant
\[
  \varphi_{ji} \;=\; f_{j} g_{i} \;:\; B_{i} \longrightarrow B_{j}
\]
et en supposant que les $f_{j}$ commutent aux produits indexés par $I$ — ce qui
est automatique si $I$ est fini, ou si les $f_{j}$ sont exacts — on obtient
\[
  \mathrm{pr}_{j}\, \varphi\bigl( (X_{i}) \bigr)
  \;=\; \prod_{i \in I} \varphi_{ji}(X_{i}) .
\]
La comonade devient donc une \emph{matrice} de foncteurs $\varphi_{ji}$, et sa
structure se réécrit en une famille de flèches
\[
  \lambda_{kji} \;:\; \varphi_{ki} \longrightarrow \varphi_{kj}\varphi_{ji} ,
\]
qui n'est autre que $f_{k} g_{i} \to f_{k}\, g_{j} f_{j}\, g_{i}$, obtenue en
insérant l'unité $\mathrm{id} \to g_{j} f_{j}$ au milieu. Les axiomes de
coünité et de coassociativité se traduisent en relations entre les
$\lambda_{kji}$, que le manuscrit renonce à écrire. Il observe seulement que
pour $i = j = k$ on retrouve la comultiplication de $\varphi_{i}$, et que les
cas où deux des trois indices coïncident sont dégénérés : ne restent que le cas
des trois indices distincts et le cas $i = k \neq j$, qui donne
\[
  \mathrm{id}_{B_{i}} \longrightarrow \varphi_{ij}\, \varphi_{ji}
  \qquad (i \neq j) .
\]

\medskip

\textbf{Deux facteurs.} Le cas $B = B' \times B''$ est traité en détail, sous
l'hypothèse que $g'$ et $g''$ sont pleinement fidèles, c'est-à-dire que
$f'g' \simeq \mathrm{id}_{B'}$ et $f''g'' \simeq \mathrm{id}_{B''}$. Il ne
reste alors, de toute la matrice, que deux foncteurs croisés
\[
  \varphi' : B'' \longrightarrow B' , \qquad
  \varphi'' : B' \longrightarrow B''
\]
et deux flèches
\[
  \lambda' : \mathrm{id}_{B'} \longrightarrow \varphi' \varphi'' , \qquad
  \lambda'' : \mathrm{id}_{B''} \longrightarrow \varphi'' \varphi' ,
\]
avec $\varphi(X', X'') = \bigl( X' \times \varphi'(X''),\ \varphi''(X') \times
X'' \bigr)$.

La catégorie $A$ s'explicite alors complètement : ses objets sont les
quadruplets $(X', X'', u', u'')$ où
\[
  u' : X' \longrightarrow \varphi'(X'') , \qquad
  u'' : X'' \longrightarrow \varphi''(X') ,
\]
soumis à la commutation des deux triangles

\begin{tikzcd}[column sep=large]
X' \arrow[r, "u'"] \arrow[rr, bend right=30, "\lambda'(X')"'] &
\varphi'(X'') \arrow[r, "\varphi'(u'')"] & \varphi'\varphi''(X')
\end{tikzcd}

\begin{tikzcd}[column sep=large]
X'' \arrow[r, "u''"] \arrow[rr, bend right=30, "\lambda''(X'')"'] &
\varphi''(X') \arrow[r, "\varphi''(u')"] & \varphi''\varphi'(X'')
\end{tikzcd}

l'axiome de coünité étant, lui, automatiquement satisfait par la forme même de
$\varphi$.\footnote{Le manuscrit dit « il semble qu'on trouve » : le calcul qui
mène à ces deux triangles est mené sur trois lignes puis biffé, et n'est pas
récrit.}

\medskip

\textbf{Le cas dégénéré.} Si $\varphi'\varphi''$ et $\varphi''\varphi'$ sont
constants de valeur l'objet final, les deux triangles sont vides de contenu et
$A$ est simplement la catégorie des quadruplets, sans condition. Les deux
foncteurs $g'$, $g''$ s'écrivent alors explicitement, et leurs images
essentielles se reconnaissent par une condition d'une simplicité frappante :
\[
  \text{image essentielle de } g' : u'' \text{ est un isomorphisme} ,
  \qquad
  \text{image essentielle de } g'' : u' \text{ est un isomorphisme} .
\]
Leur intersection est donc formée des quadruplets où $u'$ \emph{et} $u''$ sont
inversibles, et cette sous-catégorie est équivalente à la sous-catégorie pleine
de $B'$ formée des $X'$ pour lesquels $\lambda'(X')$ est un isomorphisme. La
page en tire trois conditions équivalentes, dont la teneur est que cette
intersection se réduit à un point si et seulement si $\lambda'$, ou de façon
équivalente $\lambda''$, n'est inversible que sur l'objet
trivial.\footnote{L'adjectif que la page emploie pour la catégorie
intersection est d'une lecture incertaine — la transcription le donne comme
« ponctuelle ». Le sens est en tout cas celui d'une catégorie réduite à un
objet et à son identité.}

\pagerange{13}{20}
\section*{Résolutions standard et cohomologie relative (pages 13 à 20)}

Le second texte du dossier commence page 13 et se poursuit jusqu'à la page 27,
donc au-delà de ce lot. Il applique le mécanisme précédent à l'algèbre
homologique.

\pagerange{13}{13}
\subsection*{La construction standard}

Soit $F \dashv G$ une adjonction, $F : C \to C'$ et $G : C' \to C$. On a donc
d'un côté la monade $E = GF$ sur $C$, avec son unité
$\eta : \mathrm{id}_{C} \to E$ et sa multiplication $E^{2} \to E$, et de
l'autre la comonade $E' = FG$ sur $C'$.

Itérer $E$ produit un objet cosimplicial coaugmenté dans la catégorie des
endofoncteurs de $C$ :
\[
  X \longrightarrow E(X) \rightrightarrows E^{2}(X)
  \Rrightarrow E^{3}(X) \longrightarrow \cdots
\]
dont le terme de degré $n$ est $E^{n+1}$, les cofaces sont obtenues en
insérant l'unité aux $n+2$ places possibles, et les codégénérescences en
appliquant la multiplication.\footnote{La page dit « objet simplicial » et
« augmentation $X \to \widehat{E}(X)$ » ; les flèches qu'elle écrit vont dans
le sens croissant et partent de $X$, ce qui décrit un objet \emph{cosimplicial
coaugmenté}. L'auteur se corrige lui-même page 70, où il biffe
« cosimpliciale » pour écrire « co-semi-simpliciale ». C'est la
\emph{construction standard} de Godement, la résolution bar d'aujourd'hui.} La
comonade $E'$ donne de même, sur $C'$, un objet simplicial augmenté.

Deux faits en découlent, et ce sont eux qui font l'intérêt de la construction :
pour $X'$ dans $C'$, le complexe $\widehat{E}\bigl( G(X') \bigr)$ est
contractile, et pour $X$ dans $C$, $F(X) \to F\bigl( \widehat{E}(X) \bigr)$ est
une équivalence d'homotopie. Autrement dit, la résolution devient triviale dès
qu'on la regarde du côté où elle a été construite. La raison est une
\emph{dégénérescence supplémentaire}, et le manuscrit l'énonce sous une forme
plus générale : si un foncteur $T : C \to D$ est muni d'un
$h : T \circ E \to T$ tel que $h \circ (T * \eta) = \mathrm{id}_{T}$, alors
$T\widehat{E}(X)$ est contractile pour tout $X$.

L'application visée est immédiate : si $C$ et $C'$ sont abéliennes, $F$ et $G$
additifs, et si $F$ est exact et fidèle, alors $\widehat{E}(X)$ est une
résolution de $X$. La page s'arrête au milieu de la phrase
suivante.\footnote{Elle s'interrompt sur « 2. Si $\Gamma$ est un foncteur » ;
la suite n'est pas au dossier. Le texte reprend page 14 sous un autre titre.}

\pagerange{14}{15}
\subsection*{Représentabilité : des $h_{B}$ aux $R$-modules}

Le texte reprend, sous le titre « Cohomologie relative », par une question
préalable : reconnaître, parmi tous les foncteurs, ceux de la forme
$A \mapsto \mathrm{Hom}(A, B)$.

Soit $C$ une catégorie et, pour $B$ dans $C$, soit $h_{B}$ le foncteur
$A \mapsto \mathrm{Hom}(A, B)$. Il est contravariant en $A$ et \emph{covariant}
en $B$,\footnote{La page écrit « contravariant en $B$ » ; c'est un lapsus, et
la formule qu'elle écrit à la ligne suivante — $\mathrm{Hom}(B,B') \to
\mathrm{Hom}(h_{B}, h_{B'})$, dans ce sens — dit bien la covariance.} et le
lemme de Yoneda donne
\[
  \mathrm{Hom}(B, B') \;\xrightarrow{\ \sim\ }\;
  \mathrm{Hom}\bigl( h_{B}, h_{B'} \bigr) ,
\]
la bijection réciproque s'obtenant en évaluant une transformation naturelle sur
l'identité de $B$. Le foncteur $B \mapsto h_{B}$ identifie donc $C$ à une
sous-catégorie pleine de la catégorie des foncteurs $C^{\circ} \to
\mathbf{Ab}$,\footnote{La capitale cursive qui désigne la catégorie but n'a pas
pu être lue ; le contexte — « si $C$ est additive », « homomorphismes
additifs » — impose les groupes abéliens.} additifs si $C$ l'est.

Supposons maintenant que $C$ soit une catégorie de Grothendieck : abélienne,
vérifiant AB\,3), 4), 5), et munie d'un générateur $U$. Posons
$R = \mathrm{End}_{C}(U)$ et, pour un foncteur $h$ transformant limites
inductives en limites projectives, $M(h) = h(U)$, qui est un $R$-module à
droite. Le manuscrit affirme alors que
\[
  \mathrm{Hom}(h, h') \;\xrightarrow{\ \sim\ }\;
  \mathrm{Hom}_{R^{\circ}}\bigl( M(h), M(h') \bigr) ,
\]
puis, définissant en sens inverse $M \mapsto M \otimes_{R} U$, que les deux
foncteurs sont quasi-inverses, de sorte que les foncteurs envisagés sont
« essentiellement les $h_{B}$ ».

Cet énoncé n'est correct que sous une hypothèse que la page ne formule pas.
L'énoncé vrai est le suivant. Le foncteur $\mathrm{Hom}_{C}(U, -) : C \to
\mathbf{Mod}_{R}$ est pleinement fidèle, et son adjoint à gauche $- \otimes_{R}
U$ est exact : c'est le théorème de Gabriel--Popescu, et il présente $C$ comme
une \emph{localisation} de $\mathbf{Mod}_{R}$, non comme
$\mathbf{Mod}_{R}$.\footnote{Gabriel--Popescu, 1964. Le manuscrit invoque
$\mathrm{Hom}_{C}(U, M \otimes_{R} U) \simeq M$ comme « bien connu » ; cette
formule vaut pour tout $M$ exactement quand $U$ est petit et projectif. La
page, qui ne suppose que « générateur », affirme donc une équivalence là où
seul un plongement pleinement fidèle est acquis.} Pour que ce soit une
équivalence, il faut et il suffit que $U$ soit un générateur \emph{projectif}
et petit, et l'on retombe alors sur le théorème de Morita.

C'est cette différence qui explique le mot « essentiellement » de la page, et
le « détails laissés au lecteur » qui le précède.

\pagerange{16}{17}
\subsection*{Le couple $(S,T)$ : unité et coünité}

Soient $C$ et $C'$ deux catégories additives et $S : C \to C'$ un foncteur
additif. On suppose que pour tout $X$ de $C'$ le foncteur
$A \mapsto \mathrm{Hom}_{C'}\bigl( S(A), X \bigr)$ est représentable, et de
façon fonctorielle en $X$ : on a donc un foncteur $T : C' \to C$ et une
bijection naturelle
\[
  \mathrm{Hom}_{C'}\bigl( S(A), X \bigr) \;\simeq\;
  \mathrm{Hom}_{C}\bigl( A, T(X) \bigr) .
\]
Autrement dit $S \dashv T$. La condition de représentabilité est satisfaite
notamment lorsque $C$ et $C'$ vérifient AB\,3), 4), 5), que $C$ a un
générateur, et que $S$ est exact à droite — c'est le paragraphe précédent qui
le fournit.\footnote{La marge de la page 16 porte : « il faudrait des
conditions de représentabilité ». Sous la forme donnée ici, ces conditions sont
celles du théorème d'existence d'un adjoint pour un foncteur cocontinu depuis
une catégorie de Grothendieck.}

De l'adjonction viennent l'unité et la coünité,
\[
  \eta : \mathrm{id}_{C} \longrightarrow TS , \qquad
  \varepsilon : ST \longrightarrow \mathrm{id}_{C'} ,
\]
obtenues comme d'ordinaire en évaluant la bijection sur des identités. Les deux
identités triangulaires prennent ici une lecture géométrique que la page
souligne : le composé
\[
  S(A) \xrightarrow{\;S(\eta_{A})\;} STS(A)
  \xrightarrow{\;\varepsilon_{S(A)}\;} S(A)
\]
est l'identité, ce qui fait de $S(A)$ un \emph{facteur direct} de $STS(A)$, la
rétraction étant $\varepsilon$ ; et symétriquement $T(X)$ est facteur direct de
$TST(X)$. C'est cette scission-là, et rien d'autre, qui sera le moteur de tout
ce qui suit.

Enfin, $\eta$ est un monomorphisme en chaque objet si et seulement si les
applications $\mathrm{Hom}_{C}(A,B) \to \mathrm{Hom}_{C'}\bigl( S(A), S(B)
\bigr)$ sont injectives, c'est-à-dire si et seulement si $S$ est fidèle.

\pagerange{18}{20}
\subsection*{Suites $S$-exactes et objets $S$-injectifs}

On peut maintenant refaire l'algèbre homologique en ne tenant pour exactes que
les suites que $S$ scinde. C'est l'algèbre homologique \emph{relative}.

Une suite courte $0 \to A' \to A \to A'' \to 0$ de $C$ est dite
\emph{$S$-exacte} si son image par $S$ est exacte \emph{et scindée}. Un objet
$M$ de $C$ est dit \emph{$S$-injectif} si $\mathrm{Hom}(-, M)$ transforme toute
suite $S$-exacte en une suite exacte courte de groupes abéliens.\footnote{Cette
classe de suites est ce qu'on appelle aujourd'hui une \emph{classe propre}, ou
classe admissible, de suites exactes ; l'algèbre homologique qu'on bâtit dessus
est celle de Hochschild (1956) et d'Eilenberg--Moore (1965). Le manuscrit ne
cite ni l'un ni l'autre.}

\textbf{Premier exemple.} Tout objet de la forme $T(X)$ est $S$-injectif, car
pour un tel $M$ la suite en jeu est simplement l'image par $\mathrm{Hom}(S(-),
X)$ de la suite $S$-exacte, donc exacte par scission.

\textbf{Deuxième exemple.} Supposons $S$ exact à droite et fidèle. Alors, pour
tout $A$, la suite
\[
  0 \longrightarrow A \xrightarrow{\;\eta_{A}\;} TS(A) \longrightarrow Z(A)
  \longrightarrow 0 , \qquad Z(A) = \mathrm{coker}\,\eta_{A} ,
\]
est $S$-exacte : son image par $S$ est exacte parce que $S$ est exact à droite,
et elle est scindée par la rétraction $\varepsilon_{S(A)}$ construite plus
haut.

De ces deux exemples sort la caractérisation qui organise la suite du texte :
\emph{un objet est $S$-injectif si et seulement si $\eta$ l'identifie à un
facteur direct de $TS(M)$}. En effet, s'il en est ainsi, il est facteur direct
d'un $T(X)$, donc $S$-injectif par le premier exemple ; et réciproquement, un
objet $S$-injectif rend scindée la suite $S$-exacte qui lui est attachée par le
second.\footnote{La page renvoie ici, entre crochets, « aux faits de nature
\ldots{} de la catégorie Eilenberg--Moore ». C'est la seule mention du nom dans
le dossier, et elle est juste : la description des $S$-injectifs comme facteurs
directs des $T(X)$ est le pendant, pour la classe propre, de la description des
objets libres d'une catégorie d'algèbres.}

\medskip

\textbf{Quand une suite $S$-exacte est-elle exacte ?} Si $S$ est exact et
fidèle, les deux notions coïncident. La raison est brève : $S$ étant exact,
$S\bigl( \mathrm{Ker}\,v / \mathrm{Im}\,u \bigr) = \mathrm{Ker}\,S(v) /
\mathrm{Im}\,S(u)$, et $S$ étant fidèle, le premier membre est nul si et
seulement si $\mathrm{Ker}\,v / \mathrm{Im}\,u$ l'est.\footnote{Le passage est
récrit trois fois sur les pages 19 et 20, et fortement biffé ; le haut de la
page 20 est barré en croix. Seule la version portée en interligne se lit d'un
bout à l'autre, et c'est elle qui est donnée ici.}

Le lot s'achève sur la définition de la résolution canonique. En posant
$C^{0} = TS$, puis
\[
  C^{1}(A) = C^{0}(A) / \mathrm{Im}\,A , \qquad
  C^{i}(A) = C^{0}\Bigl( C^{i-1}(A) / \mathrm{Im}\bigl( C^{i-2}(A) \bigr)
  \Bigr) \quad (i \geqslant 2) ,
\]
on obtient une $S$-résolution de $A$, fonctorielle en $A$, dont tous les termes
sont $S$-injectifs. C'est elle qui servira, dans le lot suivant, à définir les
foncteurs dérivés relatifs.
\pagerange{21}{27}
\section*{Cohomologie relative : la fin de la théorie (pages 21 à 27)}

On reprend la situation du lot précédent : deux catégories abéliennes $C$,
$C'$, un foncteur additif $S : C \to C'$ admettant un adjoint à droite $T$,
d'unité $\eta : \mathrm{id}_{C} \to TS$ et de coünité
$\varepsilon : ST \to \mathrm{id}_{C'}$. Une suite courte de $C$ est
$S$-exacte lorsque son image par $S$ est exacte et scindée ; un objet est
$S$-injectif lorsque $\mathrm{Hom}(-, M)$ transforme toute suite $S$-exacte en
suite exacte. À partir d'ici, et comme la marge de la page 23 le note une fois
pour toutes, $S$ est supposé exact et fidèle.\footnote{Le manuscrit note
$\widetilde{A}$ pour $S(A)$ et « $\sim$-exact », « $\sim$-injectif » pour
$S$-exact, $S$-injectif ; on a uniformisé.}

\pagerange{21}{21}
\subsection*{Quand les foncteurs dérivés relatifs sont les foncteurs dérivés}

Soit $\Gamma : C \to D$ un foncteur, et $R_{S}^{\bullet}\Gamma$ les foncteurs
dérivés calculés sur les $S$-résolutions. La question est de savoir quand ils
coïncident avec les $R^{\bullet}\Gamma$ ordinaires — donc, en premier lieu,
quand $R_{S}^{\bullet}\Gamma$ est bien un foncteur cohomologique, c'est-à-dire
quand une suite courte lui donne une suite exacte longue.

Trois conditions suffisent :

\begin{enumerate}
\item[(i)] $\eta_{A} : A \to C^{0}(A)$ est un monomorphisme ;
\item[(ii)] $S$ et $T$ sont exacts — donc $C^{0} = TS$ l'est, donc
$\Gamma C^{\bullet}$ aussi ;
\item[(iii)] pour toute suite courte $0 \to A' \to A \to A'' \to 0$ de $C$, la
suite
\[
  0 \longrightarrow \Gamma\bigl( C^{0}(A') \bigr) \longrightarrow
  \Gamma\bigl( C^{0}(A) \bigr) \longrightarrow
  \Gamma\bigl( C^{0}(A'') \bigr) \longrightarrow 0
\]
est exacte.
\end{enumerate}

Il suffit d'ailleurs, pour (iii), que le foncteur composé $\Gamma T$ soit
exact. Sous ces conditions $R_{S}^{\bullet}\Gamma$ est un foncteur
cohomologique et, en degré $0$, coïncide avec $R^{0}\Gamma = \Gamma$ dès que
$\Gamma$ est exact à gauche.\footnote{Une demi-page est ici récrite trois fois
puis biffée d'un trait ; seul le résultat, $R_{S}^{0}\Gamma = \Gamma =
R^{0}\Gamma$, se lit. La condition (i) est ce que la page appelle ailleurs
l'injectivité de $\eta$, équivalente à la fidélité de $S$.}

\pagerange{22}{23}
\subsection*{Les $S$-injectifs sont les facteurs directs des $T(N)$}

La caractérisation esquissée à la fin du lot précédent est ici établie.

D'un côté, tout objet $T(N)$ est $S$-injectif : appliquer $\mathrm{Hom}(-,
T(N))$ à une suite $S$-exacte revient, par adjonction, à lui appliquer
$\mathrm{Hom}\bigl( S(-), N \bigr)$, ce qui donne une suite exacte puisque la
suite image est scindée. La réciproque de cette observation est utile aussi :
si une suite courte devient exacte après $\mathrm{Hom}(-, T(N))$ pour tout $N$,
alors elle est $S$-exacte.

De l'autre côté, pour tout $X$ de $C$ la suite
\[
  \Sigma_{X} : \quad
  0 \longrightarrow X \xrightarrow{\;\eta_{X}\;} T\bigl( S(X) \bigr)
  \longrightarrow X^{*} \longrightarrow 0 ,
  \qquad X^{*} = \mathrm{coker}\,\eta_{X} ,
\]
est $S$-exacte, canoniquement scindée par la coünité.

Il suffit alors de faire $X = M$ : si $M$ est tel que $\mathrm{Hom}(-, M)$ rend
exacte toute suite de la forme $\Sigma_{X}$, alors en particulier
$\Sigma_{M}$ est scindée, donc $M$ est facteur direct de $T\bigl( S(M)
\bigr)$ ; celui-ci étant $S$-injectif, $M$ l'est aussi. D'où :

\textbf{Les objets $S$-injectifs sont exactement les facteurs directs des
objets $T(N)$.}\footnote{La page écrit « les objets injectifs sont exactement
les facteurs directs des objets $T(M)$ » ; c'est de $S$-injectivité qu'il
s'agit, la page entière étant relative. Le haut de la page 23, biffé en croix,
est une première rédaction du même argument.}

\pagerange{23}{25}
\subsection*{Morphismes effaçants, complexes droits, résolution canonique}

Trois notions se mettent en place, dont la troisième est le but.

\medskip

\textbf{$S$-résolution.} Une $S$-résolution de $A$ est une résolution
$A \to X^{\bullet}$ munie, dans le complexe augmenté image par $S$, de
scissions — c'est-à-dire de scissions des suites exactes
\[
  0 \longrightarrow S\bigl( Z^{i}(X^{\bullet}) \bigr) \longrightarrow
  S(X^{i}) \longrightarrow S\bigl( Z^{i+1}(X^{\bullet}) \bigr)
  \longrightarrow 0 \qquad (i \geqslant 0) .
\]
Ce sont ces scissions, et non le complexe seul, qui font la résolution.

\medskip

\textbf{Morphisme effaçant.} Un morphisme $\eta : X \to Y$ est dit
\emph{effaçant} si l'image de $\mathrm{Hom}\bigl( S(A), S(X) \bigr)$ dans
$\mathrm{Hom}\bigl( S(A), S(Y) \bigr)$ est toujours contenue dans l'image de
$\mathrm{Hom}(A, Y)$. Par adjonction cela revient à dire que $T(\eta) : T(X)
\to T(Y)$ se factorise par $Y$ :
\[
  T(X) \longrightarrow Y \xrightarrow{\;\eta_{Y}\;} T(Y) .
\]
La classe des morphismes effaçants est stable par composition à gauche par
n'importe quoi, et l'unité $\eta_{M} : M \to T(M)$ est toujours
effaçante.\footnote{C'est la notion d'\emph{effaçabilité} de Grothendieck, ici
transposée à la classe propre : un foncteur dérivé s'annule sur les objets que
les morphismes effaçants effacent. Le mot « effaçant » est celui de la page.}

\medskip

\textbf{Complexe droit.} Un \emph{$t$-complexe droit} sur $B$ est un complexe
$B \to Y^{\bullet}$ muni de morphismes $t : T(B) \to Y^{0}$ et
$T(Y^{i}) \to Y^{i+1}$ rendant commutatifs les triangles évidents, autrement
dit d'une factorisation de chaque différentielle à travers $T$.

Ces trois notions servent ensemble à un unique énoncé, qui est le principe
d'acyclicité de toute la théorie : \emph{une $S$-résolution s'envoie
canoniquement dans tout $t$-complexe droit, au-dessus d'un morphisme donné des
augmentations.} La construction est une récurrence, et elle est instructive :
connaissant $\bar{u}^{n} : X^{n} \to Y^{n}$, la scission $s^{n+1}$ de la
résolution fournit $X^{n+1} \to T(Y^{n})$, et la structure $t$ du but fournit
$T(Y^{n}) \to Y^{n+1}$ ; on pose
\[
  \bar{u}^{n+1} \;=\; t^{n} \circ T(\bar{u}^{n}) \circ s'^{\,n+1} .
\]
Les scissions servent à monter, la structure $t$ à redescendre.

\medskip

\textbf{La résolution canonique.} Elle s'obtient en itérant $C^{0} = TS$ sur
les conoyaux successifs :
\[
  C^{n+1}(A) \;=\; T\Bigl( S\bigl( C^{n}(A) / \mathrm{Im}\,C^{n-1}(A) \bigr)
  \Bigr) , \qquad C^{-1}(A) = A .
\]
C'est un $t$-complexe droit, c'est une $S$-résolution, et c'est un foncteur en
$A$ : le morphisme $C^{\bullet}(A) \to C^{\bullet}(B)$ attaché à $A \to B$ est
exactement celui que l'énoncé précédent produit. Les foncteurs dérivés relatifs
se calculent donc sur elle, sans choix.

\pagerange{26}{27}
\subsection*{Trois exemples}

\textbf{1. Restriction des scalaires.} Soient $U$ un anneau, $V$ un
sous-anneau, et $S : \mathbf{Mod}_{U} \to \mathbf{Mod}_{V}$ la restriction des
scalaires. Elle est exacte et fidèle. Son adjoint à droite est
$T(B) = \mathrm{Hom}_{V}(U, B)$, la coünité étant l'évaluation en $1$ :
\[
  \mathrm{Hom}_{V}\bigl( S(A), B \bigr) \;\simeq\;
  \mathrm{Hom}_{U}\bigl( A, \mathrm{Hom}_{V}(U, B) \bigr) .
\]
Les suites $S$-exactes sont celles qui se scindent sur $V$, et les $U$-modules
$S$-injectifs sont les facteurs directs des $\mathrm{Hom}_{V}(U, B)$. C'est
l'algèbre homologique relative des extensions d'anneaux, celle de
Hochschild.\footnote{Hochschild, \emph{Relative homological algebra}, 1956. Le
manuscrit ne le nomme pas.}

\textbf{2. L'autre côté.} La même restriction des scalaires admet aussi un
adjoint à \emph{gauche}, $A \mapsto U \otimes_{V} A$, et l'on se retrouve dans
la situation duale : les objets $S$-projectifs sont les facteurs directs des
$U \otimes_{V} A$, avec l'unité $a \mapsto 1 \otimes a$, qui est
$V$-linéaire, et la coünité $u \otimes a \mapsto ua$, qui est $U$-linéaire et
surjective.

\textbf{3. Le foncteur de Godement.} Soit $C$ la catégorie des faisceaux de
$\mathcal{O}$-modules sur un espace $X$, et $\widetilde{C} = \prod_{x \in X}
\mathbf{Mod}_{\mathcal{O}_{x}}$ le produit des catégories de modules sur les
fibres. Le foncteur $S$ qui à un faisceau associe la famille de ses fibres est
exact et fidèle, et son adjoint à droite est le foncteur qui à une famille
associe le faisceau des sections discontinues. Le composé $TS$ est exactement
le foncteur $C^{0}$ de Godement, et la résolution canonique ci-dessus est la
résolution flasque canonique.

C'est le troisième exemple qui explique le titre du texte, et l'ordre dans
lequel il a été écrit : la construction générale est l'abstraction de ce
cas-là.

\pagerange{28}{40}
\section*{Notes pour SGA 4, exposé IX (pages 28 à 40)}

Une chemise, page 28, porte « Remarques » et « SGA 4 IX ». Ce qui suit n'est
pas une rédaction de l'exposé mais un cahier de préparation : plan, énoncés en
sténographie, dépendances. Les feuillets ne sont d'ailleurs pas dans l'ordre de
lecture — la page 37 se poursuit page 39, et la page 38 porte la proposition
dont la page 36 tire le corollaire.

Le sujet de l'exposé est la \emph{constructibilité} : parmi les faisceaux
étales sur un schéma, reconnaître ceux qui sont localement constants à fibres
finies sur les morceaux d'une stratification finie. C'est la classe de
faisceaux sur laquelle les théorèmes de finitude de la cohomologie étale
peuvent être espérés.

\pagerange{30}{30}
\subsection*{Le plan, et les huit conditions de constructibilité}

L'exposé est prévu en trois numéros : sorite des faisceaux de torsion,
faisceaux localement constants, faisceaux constructibles. Le second numéro doit
donner les propriétés d'exactitude de la notion.

Le cœur de la page est une liste de huit conditions, présentées comme
équivalentes, dont voici celles que la lecture soutient :

\begin{enumerate}
\item[(i)] l'image inverse de $F$ sur $X^{\mathrm{red}}$ est localement
constante à fibres finies ;
\item[(ii)] il existe $f : X' \to X$ localement de présentation finie et
surjectif tel que $f^{*}(F)$ soit localement constant à fibres finies ;
\item[(iii)] $X$ étant quasi-compact et quasi-séparé, il existe une partition
finie de $X$ en localement fermés constructibles $X_{i}$ telle que
$F|_{X_{i}}$ soit localement constant à fibres finies ;
\item[(iv)] $X$ étant quasi-compact et quasi-séparé, $F$ est le noyau d'une
double flèche $F_{1} \rightrightarrows F_{0}$ entre faisceaux d'un type que la
page ne livre pas ;
\item[(v)] $X$ étant quasi-compact et quasi-séparé, une condition portant sur
les faisceaux représentables par des schémas étales de présentation finie sur
$X$ ;
\item[(vi)] $X$ étant quasi-compact et quasi-séparé, $F$ admet une filtration
finie à quotients de la forme $i_{!}(G)$, avec $i : Y \to X$ une immersion
constructible ;
\item[(vii)] une variante de (vi) faisant intervenir les $p_{*}(G_{X'})$ pour
$X'$ fini de présentation finie sur $X$ ;
\item[(viii)] pour $A = \mathbf{Z}$, faisceaux de torsion, $X$ quasi-compact
quasi-séparé : $F$ est réunion de facteurs directs de faisceaux filtrés par des
$p_{*}\bigl( i_{!}(\mathbf{Z}/n\mathbf{Z})_{U} \bigr)$, où $p : X' \to X$ est
fini, $U$ un ouvert de $X'$ et $i$ son immersion.
\end{enumerate}

Les conditions (iv), (v) et (vii) sont trop lacunaires pour être rétablies, et
on ne les a pas rétablies.\footnote{Les marques d'illisibilité y sont
littérales : la page est écrite très petit et très vite. Ce qui est donné
ci-dessus est ce que la transcription soutient, et rien de plus. La marge porte
en outre une note de rédaction : penser aussi aux faisceaux de $A$-modules et à
la $A$-constructibilité, et commencer la rédaction par la $\mathbf{Q}$- et la
$\Delta$-constructibilité.}

Suivent, en sténographie, les rubriques restantes : la définition des
constructibles de type fini ; la stabilité par limites inductives et
projectives, qui exige que $A$ soit noethérien ; la stabilité par facteurs
directs dans le cas noethérien ; la stabilité par image directe ; et la
stabilité par morphisme fini ou quasi-fini.

\pagerange{31}{31}
\subsection*{Pureté}

\textbf{Proposition.} Soit $X$ un préschéma localement noethérien régulier.
Les conditions suivantes sont équivalentes :

\begin{enumerate}
\item[(i)] pour tout sous-préschéma régulier $Y$ de $X$ et tout sous-préschéma
régulier fermé $Z$ de $Y$, purement de codimension $d$, on a
$H^{i}_{Z}(A_{Y}) = 0$ pour $i \neq 2d$ ;
\item[(ii)] la même chose pour $d = 1$ ;
\item[(iii)] la même chose lorsque $Z$ est un diviseur.
\end{enumerate}

L'implication $(i) \Rightarrow (ii)$ est triviale ; $(ii) \Rightarrow (i)$ se
fait par dévissage sur la codimension, en se ramenant au cas $d = 1$ où
$H^{2}_{Z}(A_{Y})$ est localement isomorphe à $A_{Z}$. C'est le théorème de
pureté cohomologique absolue, sous la forme d'une équivalence entre l'énoncé
général et le cas des diviseurs.\footnote{La page biffe, dans (i), la clause
« $H^{2d}_{Z}(A_{Y})$ localement isomorphe à $A_{Y}$ » : c'est
$A_{Z}$ qui est en jeu, et c'est bien $A_{Z}$ qu'elle écrit plus bas dans la
démonstration. La démonstration de $(iii) \Rightarrow (ii)$ n'est pas
lisible.}

\pagerange{32}{32}
\subsection*{La suite du plan, et le Leitfaden}

Trois propositions sont annoncées, en sténographie : écrire un faisceau
quelconque comme limite inductive, comparer un constructible à une limite
inductive de constructibles, et comparer les constructibles sur une limite
projective de préschémas (proposition e) ; caractériser la constructibilité par
la locale constance et deux autres propriétés (proposition f) ; caractériser
les faisceaux localement constants en $A$-modules (proposition g). L'exposé
doit se terminer par ce qui est propre aux courbes.

La numérotation des paragraphes est ensuite arrêtée — torseurs et faisceaux de
torsion constructibles, complexes de faisceaux, Kummer et Artin--Schreier,
dimension cohomologique des courbes algébriques — puis corrigée sur place :
une accolade et une flèche intervertissent les numéros 4 et 5.

La page porte enfin un \emph{Leitfaden}, c'est-à-dire l'arbre des dépendances
entre numéros :

\begin{tikzcd}[column sep=small, row sep=small, nodes={font=\scriptsize}]
 & 1 \arrow[d, no head] & \\
 & 2 \arrow[dl, no head] \arrow[d, no head] \arrow[dr, no head] & \\
5 & 3 \arrow[d, no head] & 6 \arrow[d, no head] \\
 & 4 & 7
\end{tikzcd}

\footnote{La page en porte deux croquis ; le premier, surchargé puis biffé, ne
se lit pas. Celui-ci est le second. Les traits sont sans pointe sur la page :
ce sont des dépendances, non des morphismes.}

\pagerange{33}{33}
\subsection*{Un critère sur les fibres géométriques}

\textbf{Énoncé.} Soient $f : X \to Y$ un morphisme, $Y$ localement noethérien,
$F$ un faisceau de $\ell$-torsion sur $Y$, et
$\xi \in H^{n}\bigl( X, f^{*}(F) \bigr)$. Pour que $\xi$ provienne de
$H^{n}(Y, F)$, il faut et il suffit que pour tout $y \in Y$ la classe induite
$\xi_{\bar{y}} \in H^{n}(X_{\bar{y}}, F_{\bar{y}})$ sur la fibre géométrique
soit nulle.

\textbf{Corollaire.} $f$ étant propre et acyclique en degrés $< n$, pour qu'il
soit \emph{globalement} acyclique pour $\ell$ il faut et il suffit que ses
fibres géométriques le soient.

La démonstration part de l'hypothèse d'acyclicité en degrés intermédiaires,
\[
  R^{q} f_{*}\bigl( f^{*}(F) \bigr) = 0 \quad (0 < q \leqslant n-1) ,
  \qquad R^{0} f_{*}\bigl( f^{*}(G) \bigr) \simeq G ,
\]
d'où, par la suite spectrale de Leray, la suite exacte
\[
  \cdots \longrightarrow H^{n}(Y, F) \longrightarrow H^{n}(X, f^{*}F)
  \longrightarrow H^{0}\bigl( Y, R^{n} f_{*}(f^{*}F) \bigr) ,
\]
qui ramène l'énoncé à ceci : une section de $R^{n} f_{*}\bigl( f^{*}(F)
\bigr)$ nulle sur chaque fibre géométrique est nulle. C'est vrai parce que les
fibres du faisceau $R^{n}f_{*}$ se calculent sur les fibres géométriques. La
page s'arrête sur « on est ramené ».\footnote{Une note de la page signale
qu'il faut supposer $n \geqslant 1$, le cas $n = 0$ demandant un énoncé
séparé.}

\pagerange{35}{37}
\subsection*{Comparer deux foncteurs cohomologiques}

C'est le lemme technique de l'exposé XII 6.2, et le seul endroit de la liasse
où l'argument est écrit en entier.

Soient $T^{\bullet}$ et $T'^{\bullet}$ deux foncteurs cohomologiques et
$\varphi^{\bullet} : T^{\bullet} \to T'^{\bullet}$ un morphisme. À une suite
exacte courte $0 \to A \to B \to C \to 0$ correspond l'échelle à cinq termes

\begin{tikzcd}[column sep=small, row sep=small, nodes={font=\scriptsize}]
T^{i-1}(B) \arrow[r] \arrow[d] & T^{i-1}(C) \arrow[r] \arrow[d] &
T^{i}(A) \arrow[r] \arrow[d] & T^{i}(B) \arrow[r] \arrow[d] &
T^{i}(C) \arrow[d] \\
T'^{i-1}(B) \arrow[r] & T'^{i-1}(C) \arrow[r] & T'^{i}(A) \arrow[r] &
T'^{i}(B) \arrow[r] & T'^{i}(C)
\end{tikzcd}

\footnote{Le quatrième terme de la ligne du haut est d'une lecture douteuse :
la transcription y lit plutôt $\varphi^{i}(B)$, mais c'est $T^{i}(B)$
qu'exige la suite exacte longue, et c'est ce qui est écrit ici.} et le lemme
des cinq, sous ses diverses formes, donne les quatre variantes que la page
énumère :

\begin{enumerate}
\item[a)] si $\varphi^{i}$ est un isomorphisme pour $i \leqslant n-1$ et un
monomorphisme pour $i = n$, et si $\varphi^{n}$ est un isomorphisme sur une
famille de cogénérateurs, alors $\varphi^{n}$ est un isomorphisme partout ;
\item[b)] si $\varphi^{i}$ est un isomorphisme pour $i \leqslant n$ et
$\varphi^{n+1}$ un monomorphisme sur une famille de cogénérateurs, alors
$\varphi^{n+1}$ est un monomorphisme partout ;
\item[c)] si les $\varphi^{i}$ sont des épimorphismes, $\varphi^{n+1}$ est un
monomorphisme dès que $T'^{n+1}$ est effaçable ;
\item[c$'$)] si les $T'^{i}$ sont effaçables pour $i > i_{0}$ et $\varphi^{i}$
épimorphe pour $i \leqslant n$, alors $\varphi^{i}$ est un isomorphisme pour
$i \leqslant n$ et un épimorphisme en degré $n+1$.
\end{enumerate}

Le point de méthode vient ensuite, et c'est lui qui vaut la peine. Grothendieck
donne un nom aux deux propriétés d'un objet $A$ qui reviennent :
\[
  I_{n}(A) : \ \varphi^{i}(A) \text{ isomorphisme pour } i \leqslant n ;
  \qquad
  J_{n}(A) : \ \varphi^{i}(A) \text{ isomorphisme pour } i < n,
  \text{ monomorphisme en } i = n .
\]
Elles s'enchaînent : $I_{n} \Rightarrow J_{n}$ ; $I_{n}$ pour tout $A$ plus
$J_{n}$ sur une famille donne $J_{n}$ pour tout $A$ ; $J_{n}$ pour tout $A$
plus $I_{n+1}$ sur une famille donne $I_{n+1}$ pour tout $A$. D'où la
conclusion : \emph{pour vérifier $I_{n}$ ou $J_{n}$ partout, il suffit de le
vérifier en degrés $\leqslant n$ sur une famille de générateurs.} C'est le
dévissage sous sa forme nue, et c'est le geste que la liasse entière
répète.\footnote{La page 37 se poursuit page 39 ; elle s'interrompt sur une
remarque, en partie illisible, disant que dans les questions de dimension
cohomologique il faut « démouler » la question et chercher des conditions pour
qu'un $T^{i}$ soit nul pour tout argument. C'est ce que la page 39 fait.}

\pagerange{36}{38}
\subsection*{Une classe de morphismes, et le lemme de Chow}

La proposition est page 38, son corollaire page 36 ; on les rétablit dans
l'ordre.

\textbf{Proposition.} Soit $M$ un ensemble de morphismes de préschémas
localement noethériens. On suppose :

\begin{enumerate}
\item[a)] l'appartenance à $M$ est locale sur la base ;
\item[b)] $M$ est stable par composition ;
\item[c)] si $g : \mathbf{P}^{r}_{Y} \to Y$ est la projection canonique,
$X' = \mathbf{P}^{1}_{Y} \times \mathbf{P}^{r-1}_{Y}$ muni du morphisme
birationnel bien connu $f : X' \to X = \mathbf{P}^{r}_{Y}$, et si $gf \in M$,
alors $g \in M$ ;
\item[d)] si $f$ est une immersion fermée, $gf$ simple et $gf \in M$, alors
$g \in M$ ; en outre $\mathbf{P}^{i}_{Y} \to Y$ et $\mathrm{id}_{Y}$ sont dans
$M$.
\end{enumerate}

Alors tout morphisme projectif et simple $f : X \to Y$ est dans
$M$.\footnote{La condition d) est portée sur la page avec une conclusion
« $gf \in M$ » qui répète l'hypothèse ; c'est $g \in M$ qu'il faut lire, sans
quoi la condition est vide et la récurrence qui suit ne dit rien.}

La démonstration est une récurrence sur $r$ : a), b), c) donnent que les
$\mathbf{P}^{r}_{Y} \to Y$ sont dans $M$, et d) descend de là à tout
sous-préschéma fermé simple de $\mathbf{P}^{r}_{Y}$, donc à tout projectif
simple.

\textbf{Corollaire.} Ajoutons la condition suivante :

\begin{enumerate}
\item[c$'$)] soient $X' \xrightarrow{f} X \xrightarrow{g} Y$ des morphismes
propres, $X_{0}$ un sous-préschéma fermé de $X$, $U = X - X_{0}$ et
$U' = f^{-1}(U)$. Si $U' \to U$ est un isomorphisme et si $X_{0} \to Y$ et
$gf : X' \to Y$ sont dans $M$, alors $g \in M$.\footnote{La page conclut
« alors $f \in M$ » ; c'est $g$ que l'argument produit et que la suite emploie
— $f$ est le morphisme birationnel de dévissage, il n'a pas à être dans $M$.}
\end{enumerate}

Alors tout morphisme \emph{propre} est dans $M$. En effet c$'$) implique c),
donc tout projectif est dans $M$ ; et le lemme de Chow, qui domine tout
morphisme propre par un morphisme projectif birationnel au-dessus d'un ouvert
dense, permet de dévisser le cas propre sur le cas projectif au moyen de
c$'$).

Cette proposition n'est pas de la géométrie : c'est un critère formel, et il
sert chaque fois qu'on veut établir une propriété pour tous les morphismes
propres sans les regarder.

\pagerange{39}{40}
\subsection*{Dévissage des faisceaux de $\ell$-torsion constructibles}

La page 39 achève la page 37 et applique le principe.

Deux énoncés d'abord, qui disent quels dévissages une annulation supporte : si
$T^{i}(B) = 0$ alors $T^{i}(A) = 0$ pour tout facteur direct $A$ de $B$ ; et si
$A$ a une filtration finie dont les quotients $A_{\alpha}$ vérifient
$T^{i}(A_{\alpha}) = 0$, alors $T^{i}(A) = 0$.

De là, pour un foncteur cohomologique sur les faisceaux de $\ell$-torsion
\emph{constructibles} sur $X$, il suffit de vérifier l'annulation sur les
faisceaux de la forme
\[
  i_{!}\Bigl( p_{*}\bigl( (\mathbf{Z}/n\mathbf{Z})_{Y'} \bigr) \Bigr)
  \;=\; q_{*}\, j_{!}\bigl( (\mathbf{Z}/n\mathbf{Z})_{Y'} \bigr) ,
\]
où $Y$ est un sous-préschéma constructible de $X$, $Y'$ un revêtement étale
fini de $Y$ de degré constant, et où le carré

\begin{tikzcd}
Y' \arrow[r, hook] \arrow[d, "p"'] & X' \arrow[d, "q"] \\
Y \arrow[r, hook, "i"'] & X
\end{tikzcd}

est celui des morphismes canoniques. Autrement dit : toute sous-catégorie
pleine des faisceaux de $\ell$-torsion constructibles, stable par facteurs
directs et par extension et contenant ces faisceaux-là, est la catégorie
entière. Et si l'on veut de plus la stabilité par sous-objet — c'est-à-dire une
sous-catégorie \emph{épaisse} — il suffit alors de vérifier que les
$q_{*}\bigl( (\mathbf{Z}/n\mathbf{Z})_{X'} \bigr)$ y sont, pour $X'$ fini sur
$X$, et l'on peut prendre $X'$ intègre, voire intègre normalisé si $X$ est
noethérien séparé.

La page 40, entièrement barrée en croix, esquisse un calcul de $R\,h^{!}$ pour
une composée $X_{0} \xrightarrow{g} Y \xrightarrow{j} X$ ; elle est trop
lacunaire pour être rendue.\footnote{Chacun des deux termes de la seule formule
lisible y est suspendu, par un signe $\wr$, à une seconde ligne illisible. On
ne l'a pas reconstituée.}
\pagerange{41}{42}
\section*{Deux feuillets sur SGA 4, exposé XII (pages 41 et 42)}

La page 41 est un crayon très pâle, illisible d'un bout à l'autre par larges
plages ; ce qu'on en tire est une remarque de rédaction, non un énoncé. Elle
dit que dans l'exposé XII les hypothèses de quasi-compacité et de
quasi-séparation servent moins qu'il n'y paraît : pour l'équivalence des
conditions de 6.5 entre elles, et dans le lemme 6.7, on peut s'en passer,
sauf pour un point de passage à la limite sur les faisceaux de
$\mathbf{Z}/n\mathbf{Z}$-modules, dont l'auteur pense d'ailleurs qu'il est
inutile lui aussi.\footnote{Le reste de la page — sur ce que devient
l'existence dans une sous-catégorie, et sur un renvoi à XIV 5.2 — n'est pas
lisible et n'a pas été reconstitué.}

La page 42 est un feuillet dactylographié de l'exposé, dont la partie (iii) est
barrée de deux traits. Au bas, de la main de l'auteur, l'instruction pour la
refaire : adapter la démonstration de XII 8.3 (iii), ou se ramener au cas où
$\pi_{0}$ est \emph{surjectif} en remplaçant $\overline{X}$ par
$\overline{X} \amalg Y$.\footnote{Le symbole que la transcription lit
$\pi_{0}$ est d'une lecture incertaine.}

\pagerange{43}{53}
\section*{Foncteurs localisants et catégories de fractions (pages 43 à 53)}

La chemise, page 43, porte le titre. Les feuillets ne sont pas dans l'ordre :
la page 44 porte les numéros 3) et 4), la page 45 les numéros 2) et 1). On les
rétablit ici dans l'ordre logique.

Le cadre est le suivant. $A$ est une catégorie abélienne, $C \subset A$ une
sous-catégorie \emph{épaisse} — pleine, et stable par sous-objets, quotients et
extensions — et $T : A \to A/C$ le foncteur de passage au quotient. Les flèches
que $T$ rend inversibles sont les \emph{$C$-isomorphismes} : celles dont le
noyau et le conoyau sont dans $C$. La question est l'existence d'un adjoint à
droite $S$ de $T$, appelé \emph{foncteur section}.

\pagerange{44}{45}
\subsection*{Le foncteur section et les objets $C$-fermés}

Trois listes de conditions équivalentes, dont la seconde est la plus utile.

\medskip

\textbf{Les objets que la localisation voit fidèlement.} Soit $T \dashv S$ une
adjonction quelconque et $x$ un objet de $A$. Si $S$ est pleinement fidèle, les
conditions suivantes sont équivalentes :

\begin{enumerate}
\item[a)] $x$ est isomorphe à un objet de la forme $S(y)$ ;
\item[b)] l'unité $x \to ST(x)$ est un isomorphisme ;
\item[c)] pour tout $x'$, l'application
$\mathrm{Hom}(x', x) \to \mathrm{Hom}(Tx', Tx)$ est bijective.
\end{enumerate}

L'implication b) $\Rightarrow$ a) ne demande rien de $S$ ; c'est la réciproque
qui a besoin de la pleine fidélité. On dit alors que $x$ est \emph{$C$-fermé},
ou \emph{local}, ou \emph{libre au-dessus de $Tx$} — les trois mots sont dans
le dossier.

\medskip

\textbf{Quand $S$ est-il pleinement fidèle ?} Pour une adjonction $T \dashv S$
quelconque, il y a équivalence entre :

\begin{enumerate}
\item[a)] $S$ est pleinement fidèle ;
\item[b)] la coünité $TS \to \mathrm{id}$ est un isomorphisme ;
\item[c)] $T$ est un foncteur de passage au quotient, c'est-à-dire : en notant
$M$ l'ensemble des flèches de $A$ que $T$ rend inversibles, le foncteur induit
$A[M^{-1}] \to B$ est une équivalence.
\end{enumerate}

C'est la caractérisation d'une \emph{localisation réflexive} : une adjonction
dont la coünité est inversible, autrement dit une monade idempotente
$\Lambda = ST$.\footnote{Le manuscrit ne prononce pas le mot ; il écrit
$Lx$ pour $ST x$ page 45, et $\Lambda$ page 46.}

\medskip

\textbf{L'existence du foncteur section.} Enfin, dans le cas $B = A/C$, les
conditions suivantes sur $C$ sont équivalentes :

\begin{enumerate}
\item[a)] $T$ admet un adjoint à droite $S$ ;
\item[b)] tout $x$ de $A$ a un plus grand sous-objet appartenant à $C$, et si
ce sous-objet est nul, $x$ se plonge dans un objet $C$-fermé ;
\item[c)] tout $x$ de $A$ admet un $C$-isomorphisme $x \to \widetilde{x}$ vers
un objet $C$-fermé.
\end{enumerate}

Ce $\widetilde{x}$ est alors unique à isomorphisme unique près, et n'est autre
que $ST(x)$.

\pagerange{46}{47}
\subsection*{Deux sous-catégories localisantes emboîtées}

On se donne maintenant $C$ une catégorie de Grothendieck — abélienne, à limites
inductives exactes, avec générateurs — et deux sous-catégories
$C'' \subset C' \subset C$ \emph{localisantes}, c'est-à-dire épaisses et
stables par limites inductives.

À chacune est attaché un foncteur de torsion :
\[
  \Gamma'(F) = \text{le plus grand sous-objet de } F \text{ appartenant à }
  C' , \qquad \Gamma''(F) = \text{de même pour } C'' ,
\]
et l'on pose $(\Gamma'/\Gamma'')(F) = \Gamma'(F)/\Gamma''(F)$. Ces foncteurs se
définissent aussi comme adjoints à droite des inclusions, ce que la marge de la
page 46 note : $\Gamma' = i'^{\,!}$, $\Gamma'' = i''^{\,!}$.\footnote{C'est la
notation que le manuscrit se recommande à lui-même en marge, et c'est celle
qu'on emploiera plus bas, page 48. Elle a l'avantage de rendre l'adjonction
visible : $\mathrm{Hom}_{C}\bigl( i'_{*}(x'), x \bigr) \simeq
\mathrm{Hom}_{C'}\bigl( x', i'^{\,!}(x) \bigr)$.}

D'autre part, la localisation de $C'$ le long de $C''$ donne
$T : C' \to C'/C''$ avec son foncteur section $S$, d'où
$\Lambda = ST : C' \to C'$ et l'unité $\mathrm{id}_{C'} \to \Lambda$. Cette
dernière est caractérisée par deux propriétés : $x' \to \Lambda(x')$ est un
$C''$-isomorphisme, et $\Lambda(x')$ est $C''$-fermé.

\textbf{Lemme.} $\Gamma'$, $\Gamma''$ et $\Lambda$ existent, et $\Gamma'$
commute aux sous-objets : si $x \hookrightarrow y$, le plus grand sous-objet de
$x$ dans $C''$ est $x \cap y''$, où $y''$ est celui de $y$. D'où un
monomorphisme naturel
\[
  \Gamma'(x)/\Gamma''(x) \hookrightarrow
  R^{0}(\Gamma'/\Gamma'')(x) .
\]

Vient alors la question, et elle est intéressante parce que l'auteur y répond
lui-même, dans la marge, par la négative. Le morphisme naturel
\[
  R^{0}(\Gamma'/\Gamma'')(x) \longrightarrow \Lambda\,\Gamma'(x)
\]
— qui existe parce que $\Lambda\Gamma'$ est exact à gauche — est-il un
isomorphisme ? \textbf{Non, pas en général.} Il l'est lorsque les objets
injectifs de $C''$ sont encore injectifs dans $C$.\footnote{La page porte
« \textbf{Prop.} C'est là un isom.\ (?) », et en marge, d'une autre encre :
« Faux en général ! OK si les injectifs de $C''$ sont injectifs dans $C$. »
On a écrit l'énoncé corrigé. C'est exactement l'hypothèse qui reparaît page 48
sous le numéro (i), et le lien n'est pas fortuit : c'est la même condition qui
gouverne tout ce qui suit.}

Sous cette hypothèse, on obtient : $R^{0}(\Gamma'/\Gamma'')(x)$ est
$C''$-fermé, et $\Gamma'(x)/\Gamma''(x) \to R^{0}(\Gamma'/\Gamma'')(x)$ est un
$C''$-isomorphisme. Et sans hypothèse aucune, si $x$ est injectif, alors
$\Gamma'(x)/\Gamma''(x)$ est déjà $C''$-fermé.

\pagerange{48}{50}
\subsection*{Quand le quotient conserve les injectifs}

Cinq conditions sont mises en regard, pour une sous-catégorie localisante $C$
d'une catégorie abélienne $A$, avec $i_{*} : C \to A$ l'inclusion et
$i^{!}$ son adjoint à droite :

\begin{enumerate}
\item[(i)] $i_{*}$ transforme les injectifs en injectifs ;
\item[(ii)] pour tout injectif $x$ de $A$, le plus grand sous-objet de $x$
appartenant à $C$ est encore injectif dans $A$ ;
\item[(iii)] pour tout objet de $C$, une enveloppe injective dans $A$ est dans
$C$ — autrement dit $C$ est stable par enveloppes injectives ;
\item[(iv)] le morphisme naturel $R^{0}\bigl( \mathrm{id}/i_{*}i^{!} \bigr) \to
ST$ est un isomorphisme, c'est-à-dire : pour $x$ injectif dans $A$, le quotient
$x/i_{*}i^{!}(x)$ est $C$-fermé ;
\item[(v)] $T : A \to A/C$ transforme les injectifs en injectifs.
\end{enumerate}

Le lemme qui les relie est celui de la page 50 : \emph{si $x$ est dans $C$ et
$\overline{x}$ une enveloppe injective de $x$ dans $A$, alors
$i^{!}(\overline{x})$ est une enveloppe injective de $x$ dans $C$} ; d'où
que tout injectif de $C$ est, à isomorphisme près, de la forme $i^{!}(y)$ avec
$y$ injectif dans $A$.

De là (i) $\Leftrightarrow$ (ii) $\Leftrightarrow$ (iii). La première
implication tient à ce que $i^{!}$, adjoint à droite d'un foncteur exact,
conserve les injectifs. La seconde se lit sur l'enveloppe : $i_{*}i^{!}
(\overline{x})$ est un objet de $C$ contenant $x$ et extension essentielle de
$x$, donc son enveloppe injective dans $C$ ; si (ii) vaut, il est injectif dans
$A$, ce qui est (iii). Enfin (ii) $\Rightarrow$ (iv), parce que
$i_{*}i^{!}(x)$ est $C$-fermé dans $x$ lorsque $x$ est
injectif.\footnote{Le manuscrit écrit, dans (iv), « $x/x'$ est $C''$-fermé »,
avec le double prime des pages précédentes, alors que la page n'a introduit que
$C$ ; on a écrit $C$-fermé. Les implications (iv) $\Rightarrow$ (v) et le
retour ne sont pas sur les pages conservées.}

\pagerange{52}{53}
\subsection*{L'existence du foncteur section, d'après Gabriel}

Les pages 52 et 53 reprennent la liste des conditions de $C$-fermeture, dans le
cadre général d'une catégorie de fractions, et en donnent une cinquième forme :
$x$ est dans l'image essentielle de $S$, c'est-à-dire isomorphe à un objet
$S(z)$.\footnote{La page écrit ici « isomorphe à un objet de la forme $Tz$ » ;
c'est $Sz$, l'image essentielle de $S$ étant en jeu, et c'est ce que la ligne
précédente annonce.} Dans le cas abélien, la condition prend la forme
concrète : tout sous-objet de $x$ appartenant à $C$ est nul, et $x$ est facteur
direct dans tout $y \supset x$ tel que $y/x$ soit dans $C$.

\textbf{Corollaire 1.} Pour que l'adjoint $S$ existe, il faut et il suffit que
tout $x$ de $A$ admette une flèche $x \to x'$ telle que $Tx \to Tx'$ soit un
isomorphisme et que $x'$ soit $C$-fermé.

\textbf{Corollaire 2 (Gabriel).} Dans le cas $A$ abélienne, $B = A/C$ avec $C$
épaisse, $S$ existe si et seulement si :

\begin{enumerate}
\item[a)] tout $x$ de $A$ admet un plus grand sous-objet appartenant à
$C$ ;\footnote{La page écrit « $\forall x \in \mathrm{Ob}\,C$ » ; c'est
$\mathrm{Ob}\,A$ qu'il faut lire, la condition étant vide sinon.}
\item[b)] si ce sous-objet est nul, $x$ se plonge dans un objet $C$-fermé.
\end{enumerate}

Et lorsque $C$ est stable par limites inductives — la condition a) étant alors
automatique — le foncteur $T$ transforme les injectifs en
injectifs.\footnote{C'est le corollaire 3 de Gabriel, auquel les pages 47, 48
et 50 renvoient. Gabriel, \emph{Des catégories abéliennes}, 1962. La condition
a) est ce qu'on appelle aujourd'hui l'existence d'une théorie de torsion
héréditaire attachée à $C$.}

\pagerange{54}{60}
\section*{Le théorème de recollement d'Artin et Cartier (pages 54 à 60)}

La chemise, page 54, porte : « Théorèmes de descente de M. Artin --
P. Cartier », et un pointage d'archiviste de onze pages, qui déborde ce lot :
le texte se poursuit jusqu'à la page 68.

\medskip

\noindent\textbf{Une convention, à lire avant les formules.} Le manuscrit
prend, pour son exemple directeur, un espace $X$, un ouvert $U$ et le fermé
complémentaire $Y = X - U$, et il nomme
\[
  i \;\longleftrightarrow\; \text{l'ouvert } U , \qquad
  j \;\longleftrightarrow\; \text{le fermé } Y .
\]
\emph{C'est l'inverse exact de l'usage d'aujourd'hui}, où $j$ désigne l'ouvert
et $i$ le fermé. On a gardé ses lettres, pour que cette lecture puisse être
posée à côté de la transcription formule par formule ; le lecteur qui vient de
la littérature moderne doit donc échanger les deux lettres, et non pas les
lire.\footnote{Le choix est délibéré et il n'y en a pas de bon. Renommer aurait
donné des formules d'apparence familière et de sens inverse, ce qui est le pire
des deux dangers. La cohérence interne du manuscrit, elle, est parfaite : avec
son $j$ fermé, $j_{!} = j_{*}$, et avec son $i$ ouvert, $i^{*} = i^{!}$, ce que
la page 58 écrit en toutes lettres.}

\pagerange{56}{57}
\subsection*{Les quatre conditions et la catégorie des triples}

Soient $C$, $C'$, $C''$ trois catégories et
\[
  f^{*} = (i^{*}, j^{*}) : C \longrightarrow C' \times C'' = D
\]
un foncteur. On suppose que $i^{*}$ et $j^{*}$ ont des adjoints à droite
$i_{*}$, $j_{*}$ ; alors $f^{*}$ en a un aussi, à savoir
$f_{*}(G,H) = i_{*}(G) \times j_{*}(H)$, comme le montre le calcul immédiat
\[
  \mathrm{Hom}\bigl( F, f_{*}(G,H) \bigr) \simeq
  \mathrm{Hom}(i^{*}F, G) \times \mathrm{Hom}(j^{*}F, H) \simeq
  \mathrm{Hom}\bigl( f^{*}F, (G,H) \bigr) .
\]

On pose $\varphi = j^{*} i_{*} : C' \to C''$ — le foncteur qui, d'un objet sur
l'ouvert, produit son « bord » sur le fermé — et l'on suppose de plus :

\begin{enumerate}
\item[(a)] $i_{*}$ et $j_{*}$ sont pleinement fidèles, c'est-à-dire
$i^{*}i_{*} \simeq \mathrm{id}$ et $j^{*}j_{*} \simeq \mathrm{id}$ ;
\item[(b)] $i^{*}j_{*}(H)$ est l'objet final de $C'$ pour tout $H$ — ce qui
dit que ce qui vit sur le fermé ne se voit pas sur l'ouvert ;
\item[(c)] les noyaux et conoyaux existent dans $C$, et $f^{*}$ y commute ;
\item[(d)] $f^{*}$ est conservatif : si $i^{*}\lambda$ et $j^{*}\lambda$ sont
des isomorphismes, $\lambda$ en est un.
\end{enumerate}

\textbf{Théorème (M. Artin -- P. Cartier).} Sous ces conditions, le foncteur
naturel
\[
  C \longrightarrow K
\]
est une \emph{équivalence de catégories}, où $K$ est la catégorie des triples
$(G, H, u)$ avec $G$ dans $C'$, $H$ dans $C''$ et
$u \in \mathrm{Hom}\bigl( H, \varphi(G) \bigr)$.

C'est le même mécanisme que dans tout le dossier : $f^{*}$ perd de
l'information, $\varphi$ mesure ce qui a été perdu, et un objet de $C$ est un
objet de $C' \times C''$ muni d'une donnée de recollement. La condition (b) est
ce qui rend la donnée de recollement aussi simple qu'une seule flèche : elle
force $f^{*}f_{*}(G,H) \simeq \bigl( G, \varphi(G) \times H \bigr)$, et les
recollements possibles sont alors paramétrés par un $u : H \to \varphi(G)$
arbitraire.\footnote{Dans le cas abélien, où la page se placera dès la page 58,
l'objet final est $0$ et la condition (b) s'écrit $i^{*}j_{*} = 0$. La page 57
donne l'exemple qui a tout dicté : $C$ les faisceaux sur $X$, $C'$ les
faisceaux sur l'ouvert $U$, $C''$ les faisceaux sur le fermé $Y$. Le nom
\emph{recollement}, et l'axiomatisation qu'on lui donne aujourd'hui, sont de
Beĭlinson, Bernstein et Deligne, 1982 ; le manuscrit ne dispose que du nom de
« théorème de descente » et de l'attribution qu'il porte en chemise.}

La page ajoute une condition d$'$) équivalente à (d), et c'est la forme sous
laquelle on énonce d'ordinaire le recollement : pour tout $F$ de $C$, le carré

\begin{tikzcd}
F \arrow[r] \arrow[d] & i_{*} i^{*} F \arrow[d] \\
j_{*} j^{*} F \arrow[r] & j_{*}\bigl( \varphi(i^{*}F) \bigr)
\end{tikzcd}

est cartésien. Un objet est donc bien le produit fibré de ses deux
restrictions au-dessus de leur comparaison.

\pagerange{58}{58}
\subsection*{Les six foncteurs}

On passe ici aux catégories abéliennes et aux foncteurs additifs, et le tableau
se complète : chacun des deux foncteurs $i^{*}$ et $j_{*}$ acquiert un adjoint
du côté qui lui manquait.

\begin{tikzcd}[column sep=huge]
C'' \arrow[r, "j_{*}" description] &
C \arrow[l, bend right=45, "j^{*}"'] \arrow[l, bend left=45, "j^{!}"]
  \arrow[r, "i^{*}" description] &
C' \arrow[l, bend right=45, "i_{!}"'] \arrow[l, bend left=45, "i_{*}"]
\end{tikzcd}

avec les quatre adjonctions
\[
  i_{!} \dashv i^{*} \dashv i_{*} , \qquad j^{*} \dashv j_{*} \dashv j^{!} .
\]
Les exactitudes sont celles qu'on attend : $j_{!} = j_{*}$, $j^{*}$,
$i^{*} = i^{!}$ et $i_{!}$ sont exacts ; $j^{!}$ et $i_{*}$ sont seulement
exacts à gauche, comme adjoints à droite. Le foncteur de comparaison
$\varphi = j^{*}i_{*} : C' \to C''$ est exact à gauche.

Les quatre annulations croisées
\[
  i^{*} j_{*} = 0 , \qquad j^{!} i_{*} = 0 , \qquad
  j^{*} i_{!} = 0 , \qquad j^{!} i_{!} = 0
\]
disent que l'ouvert et le fermé ne se voient pas l'un l'autre. Et les quatre
composés $C \to C$ portent leurs noms géométriques :
\[
  j_{*}j^{!}F = \Gamma_{Y}F , \quad j_{*}j^{*}F = F^{Y} , \quad
  i_{*}i^{*}F = \Gamma_{U}F , \quad i_{!}i^{*}F = F^{U} ,
\]
soit les sections à support dans le fermé, la restriction au fermé prolongée,
les sections sur l'ouvert, et le prolongement par zéro depuis l'ouvert.

\pagerange{58}{60}
\subsection*{Quand $j^{*}$ est exact}

C'est la question sur laquelle le texte s'installe, et elle occupe le reste du
lot.

\textbf{Proposition 1.} Partons de $C'' \xrightarrow{j_{*}} C
\xrightarrow{i^{*}} C'$, avec $j_{*}$ pleinement fidèle, $i^{*}$ un passage au
quotient exact — donc $i^{*}j_{*} = 0$ — et $j_{*} \simeq \mathrm{Ker}\,i^{*}$.
Alors :

\begin{enumerate}
\item[(i)] si $i_{!}$ existe, alors $j^{*}$ existe, on a la suite exacte
\[
  i_{!} i^{*} F \longrightarrow F \longrightarrow j_{*} j^{*} F
  \longrightarrow 0 ,
\]
et $j^{*}j_{*} \simeq \mathrm{id}_{C''}$, $i^{*}i_{!} \simeq \mathrm{id}_{C'}$,
$j^{*}i_{!} = 0$ ;
\item[(i bis)] dualement, si $i_{*}$ existe, alors $j^{!}$ existe, on a
\[
  0 \longrightarrow j_{*} j^{!} F \longrightarrow F \longrightarrow
  i_{*} i^{*} F ,
\]
et $j^{!}j_{*} \simeq \mathrm{id}_{C''}$, $i^{*}i_{*} \simeq \mathrm{id}_{C'}$,
$j^{!}i_{*} = 0$.
\end{enumerate}

Noter qu'en général on ne peut pas compléter ces suites par des zéros aux
extrémités manquantes : c'est précisément l'objet du corollaire qui
suit.\footnote{La marge porte une question restée sans réponse sur la page :
si $j^{*}$ existe, $i_{!}$ existe-t-il ? Elle est reprise page 62, en anglais,
et laissée ouverte là aussi.}

\medskip

\textbf{Corollaire 1.} Sous les conditions de la proposition 1, et en supposant
$i_{!}$ existant, les conditions suivantes sont équivalentes :

\begin{enumerate}
\item[a)] $j^{*}$ est exact et $i_{!} \simeq \mathrm{Ker}\,j^{*}$ ;
\item[b)] $i_{!}$ a une image essentielle épaisse ;
\item[c)] $i_{!}i^{*}(F) \to F$ est un monomorphisme pour tout $F$,
c'est-à-dire la suite
\[
  0 \longrightarrow \underbrace{i_{!} i^{*} F}_{F^{U}} \longrightarrow F
  \longrightarrow \underbrace{j_{*} j^{*} F}_{F^{Y}} \longrightarrow 0
\]
est exacte ;
\item[d)] $j^{*}$ est exact.
\end{enumerate}

La démonstration est un cycle. a) $\Rightarrow$ b) est immédiat. Pour
b) $\Rightarrow$ c), soit $N$ le noyau de $i_{!}i^{*}F \to F$ ; appliquant
$i^{*}$, qui est exact, et utilisant $i^{*}i_{!}i^{*}F \simeq i^{*}F$, on
trouve $i^{*}(N) = 0$ ; mais b) donne $N = i_{!}(M)$, et
$i^{*}i_{!}(M) = 0$ entraîne $M = 0$ puisque $i^{*}$ est un passage au
quotient, donc $N = 0$.

Pour d) $\Rightarrow$ b), on empile trois copies de la suite de (i) pour une
suite exacte $F \to F' \to F''$ et l'on tire, de l'injectivité de
$i_{!}i^{*}F'' \to F''$, que $j_{*}j^{*}F \to j_{*}j^{*}F'$ est injectif, donc
$j^{*}$ exact — il l'était déjà à droite. Puis on identifie l'image de $i_{!}$
au noyau de $j^{*}$ : si $F = i_{!}G$, alors $i_{!}i^{*}F \to F$ est un
isomorphisme, donc $j_{*}j^{*}F = 0$ ; et réciproquement $j^{*}F = 0$ force
$i_{!}i^{*}F \simeq F$.

Enfin d) $\Rightarrow$ c) : avec $N$ comme ci-dessus, $j^{*}N \to
j^{*}i_{!}i^{*}F = 0$ donne $j^{*}N = 0$, et $i^{*}N = 0$ comme plus haut ; le
premier force $N = j_{*}M$ et le second $M = 0$.

\medskip

\textbf{Deux remarques, et une symétrie qui va se briser.} Les conditions
ci-dessus ne résultent pas de la seule exactitude de $i_{!}$ : la localisation
des modules en fournit le contre-exemple.\footnote{Le contre-exemple est écrit
page 65, dans le lot suivant : $A$ un anneau de valuation discrète, $S =
\{\pi\}$, $C''$ les modules tués par une puissance de $\pi$, $C' = C_{S^{-1}A}$
et $i^{*}$ la localisation.} Et la situation duale n'est pas vérifiée dans les
cas visés : $j^{!}$ ne sera pas exact, l'image de $i_{*}$ ne sera pas épaisse,
$F \to i_{*}i^{*}F$ ne sera pas surjectif.

Sous les conditions équivalentes du corollaire, en revanche, il y a une
symétrie remarquable entre $C'$ et $C''^{\circ}$. La situation réduite

\begin{tikzcd}[column sep=huge]
C'' \arrow[r, "j_{*}" description] &
C \arrow[l, bend right=35, "j^{*}"'] \arrow[r, "i^{*}" description] &
C' \arrow[l, bend right=35, "i_{!}"']
\end{tikzcd}

vérifie en effet : $i^{*}$ est un passage au quotient exact de noyau $j_{*}$ ;
$j^{*}$ est un passage au quotient exact de noyau $i_{!}$ ; $j^{*} \dashv
j_{*}$ et $i_{!} \dashv i^{*}$.\footnote{La page écrit, pour la première de ces
trois lignes, « $i^{*}$ passage au quotient exact de noyau $i_{!}$ ». C'est
$j_{*}$ : le noyau de la restriction à l'ouvert est formé de ce qui est porté
par le fermé. La transcription signale déjà que la symétrie annoncée l'exige.}
De façon précise, la situation obtenue en passant aux catégories opposées,

\begin{tikzcd}[column sep=huge]
C'^{\circ} \arrow[r, "i_{!}^{\circ}" description] &
C^{\circ} \arrow[l, bend right=35, "i^{*\circ}"'] \arrow[r, "j^{*\circ}" description] &
C''^{\circ} \arrow[l, bend right=35, "j_{*}^{\circ}"']
\end{tikzcd}

satisfait aux hypothèses de départ. La dernière phrase du lot est un
avertissement : « Cette symétrie va être détruite dans ce qui suit. » Le lot
suivant le vérifie.
\pagerange{61}{68}
\section*{Recollement : la suite (pages 61 à 68)}

On garde la convention du lot précédent : dans ce texte, $i$ est l'immersion de
l'\emph{ouvert} et $j$ celle du \emph{fermé}, ce qui est l'inverse de l'usage
d'aujourd'hui. L'auteur numérote lui-même ces pages 5, 5\,bis, 5\,ter, 6, dans
la suite commencée page 56, et il passe à l'anglais de la page 62 à la page
67.\footnote{Le changement de langue est sans signification apparente : il
survient au milieu d'un raisonnement et cesse de même. La transcription garde
l'anglais ; cette lecture-ci est en français d'un bout à l'autre, comme le
reste du dossier.}

\pagerange{61}{63}
\subsection*{Le corollaire 2, et sa réciproque}

\textbf{Corollaire 2.} Sous les conditions de la proposition 1, les conditions
suivantes sont équivalentes :

\begin{enumerate}
\item[a)] $j^{*}$ existe et est exact ;
\item[b)] $i_{!}$ existe et a une image essentielle épaisse ;
\item[b$'$)] $i_{!}$ existe et $i_{!}i^{*}F \to F$ est un monomorphisme pour
tout $F$.
\end{enumerate}

Dans ce cas $i_{!} = \mathrm{Ker}\,j^{*}$, $j^{*} = \mathrm{Coker}\,i_{!}$, et
$i_{!}$ est exact. La réciproque est fausse, et la page le souligne : il ne
\emph{suffit} pas que $j^{*}$ et $i_{!}$ existent avec $i_{!}$
exact.\footnote{Une phrase insérée entre les lignes, commençant par « néc. »,
n'a pas pu être lue ; une autre, « néc.\ pl.\ fid. », est portée sous la
condition a).}

L'implication qui reste, b) $\Rightarrow$ a), est établie page 62 par une
construction directe, et elle mérite d'être vue parce qu'elle produit $i_{!}$
au lieu de le supposer. On définit
\[
  \tau(F) = \mathrm{Ker}\bigl( F \longrightarrow j_{*}j^{*}F \bigr) .
\]
Comme $j^{*}$ est exact et $j_{*}$ pleinement fidèle, $\tau$ est un foncteur
exact, nul sur les objets de la forme $j_{*}H$. Il se factorise donc par
$i^{*}$, et l'on obtient
\[
  \tau(F) \simeq i_{!}\bigl( i^{*}F \bigr)
\]
avec un foncteur bien défini $i_{!} : C' \to C$.\footnote{La page écrit
$i_{!} : C'' \to C$, ici comme aux pages suivantes ; c'est $C' \to C$, $i_{!}$
étant adjoint à gauche de $i^{*} : C \to C'$. La transcription porte déjà le
« sic ».} Reste à vérifier l'adjonction $\mathrm{Hom}(i_{!}G, F) \simeq
\mathrm{Hom}(G, i^{*}F)$, ce qui, $j_{*}j^{*}F$ étant le plus grand quotient de
$F$ à support dans le fermé et $\mathrm{Im}\,j_{*} = \mathrm{Ker}\,i^{*}$,
résulte de la construction même de la catégorie quotient.

Reste la question, posée deux fois dans le dossier et laissée ouverte : l'existence
de $j^{*}$, exact ou non, entraîne-t-elle celle de $i_{!}$ ?

\pagerange{63}{66}
\subsection*{Une deuxième généralisation, et le contre-exemple}

La page 63 dresse le tableau des annulations et des identifications :
\[
  \begin{array}{lll}
    j^{*} i_{!} = 0 & i_{!} = \mathrm{Ker}\,j^{*} &
      j^{*} = \mathrm{Coker}\,i_{!} \\[2pt]
    i^{*} j_{*} = 0 & j_{*} = \mathrm{Ker}\,i^{*} &
      i^{*} = \mathrm{Coker}\,j_{*}
  \end{array}
\]
puis pose une deuxième généralisation, portant cette fois sur la paire
$(i^{*}, i_{!})$ plutôt que sur $(j_{*}, j^{*})$, avec trois variantes
formellement distinctes de la condition d'épaisseur : image épaisse,
$i_{!}i^{*}F \to F$ injectif, $i_{!} = \mathrm{Ker}\,j^{*}$. La situation
entière, note-t-il, est contenue dans la seule suite exacte
\[
  0 \longrightarrow i_{!} i^{*} F \longrightarrow F \longrightarrow
  j_{*} j^{*} F \longrightarrow 0 .
\]

La page 64 énumère, dans le cas où $i_{*}$ existe aussi, trois conditions qui
s'entraînent : $j^{!}$ exact, $i_{*}$ exact, $\varphi = j^{*}i_{*}$
exact.\footnote{Les implications sont portées par des flèches doubles dont le
sens n'est pas toujours net sur la page ; on les a laissées telles, comme une
chaîne. La page porte en outre, sous b), deux conséquences : l'épaisseur d'un
noyau et la surjectivité de $F \to i_{*}i^{*}F$.}

\medskip

\textbf{Le contre-exemple.} La page 65 pose la question qui achève la symétrie
du lot précédent : suffit-il que $i_{!}$ soit exact ? La réponse est
\emph{non}, et le contre-exemple est le suivant. Soit $A$ un anneau de
valuation discrète d'uniformisante $\pi$, et $S = \{\pi^{n}\}$. Prenons pour
$C$ les $A$-modules, pour $C''$ ceux qui sont tués par une puissance de $\pi$,
pour $C'$ les $S^{-1}A$-modules, pour $i^{*}$ la localisation et pour $i_{*}$
la restriction des scalaires. Alors $i_{*}$ est exact, mais l'adjoint de
l'autre côté ne l'est pas.

C'est la vérification annoncée à la dernière ligne du lot précédent : la
symétrie entre l'ouvert et le fermé, valable au niveau des axiomes, se brise
dès qu'on demande de l'exactitude.

\medskip

\textbf{Les quatre types de situation.} La page 66 les recense, et c'est le
résumé le plus net de tout ce texte :

\begin{itemize}
\item un foncteur exact $i^{*}$, qui est un passage au quotient, et qui a deux
adjoints $i_{!}$, $i_{*}$ dont l'un est exact et pleinement fidèle à image
épaisse ;
\item un foncteur exact $j_{*}$, pleinement fidèle à image épaisse, qui a deux
adjoints $j^{*}$, $j^{!}$ dont l'un, $j^{*}$, est exact et est un passage au
quotient ;
\item les situations mixtes, mêlant $i^{*}$, $i_{*}$, $j^{*}$, $j_{*}$ ;
\item les foncteurs seulement exacts à gauche.
\end{itemize}

Et l'énoncé qui les relie : partant de $i_{*}$ pleinement fidèle et $i^{*}$
passage au quotient, il y a équivalence entre — a) $i^{*}$ a un second adjoint
$i_{!}$ à image épaisse ; b) $i^{*}$ est exact, du type $C \to C/j_{*}C''$
pour une sous-catégorie épaisse $C''$, et $j^{*}$ existe et est exact ;
c) la situation est celle, « standard », que définit un foncteur exact
$\varphi$. La construction de $j^{*}$ à partir de a) est immédiate :
\[
  j^{*}F = \mathrm{Coker}\bigl( i_{!} i^{*} F \longrightarrow F \bigr) .
\]

\pagerange{67}{68}
\subsection*{Les douze foncteurs, et les limites}

La page 67 fait le compte. Dans la situation

\begin{tikzcd}[column sep=huge]
C'' \arrow[r, bend left=25, "j_{*}"] &
C \arrow[l, bend left=25, "j^{*}"] \arrow[r, bend left=25, "i^{*}"] &
C' \arrow[l, bend left=25, "i_{!}"]
\end{tikzcd}

on dispose de deux suites exactes,
\[
  0 \longrightarrow F_{U} \longrightarrow F \longrightarrow F^{Y}
  \longrightarrow 0 , \qquad
  0 \longrightarrow \Gamma_{Y}F \longrightarrow F \longrightarrow \Gamma_{U}F ,
\]
la première bâtie sur les $f_{!}f^{*}$, la seconde sur les $f_{*}f^{!}$, pour
$f = i$ et $f = j$. Et de douze foncteurs en tout : les six de base
$j_{*}, j^{*}, j^{!}$ et $i^{*}, i_{!}, i_{*}$ ; les quatre composés
$C \to C$, soit $i_{!}i^{*}$, $i_{*}i^{*}$, $j_{*}j^{*}$, $j_{*}j^{!}$ ; le
foncteur de comparaison $\varphi = j^{*}i_{*} : C' \to C''$ ; et son prolongé
$j_{*}\varphi : C' \to C$.

La page 68 revient au français et clôt sur les limites projectives. La paire
$(i^{*}, j^{*})$ y commute, et il suffit, pour que
$\varprojlim_{k}(G_{k}, H_{k}, u_{k})$ existe dans la catégorie des triples,
que $\varprojlim G_{k}$ et $\varprojlim H_{k}$ existent dans $C'$ et $C''$. Le
mot de la fin est une comparaison : pour les propriétés d'existence, la paire
$(i^{*}, j^{*})$ est « aussi bonne que » le seul foncteur
$\varphi : C' \to C''$ — si $\varphi$ commute aux limites projectives finies,
ou quelconques, et si ces limites existent dans $C'$ et $C''$, alors elles
existent dans $C$ et la paire y commute.

\pagerange{70}{70}
\section*{L'objet co-semi-simplicial d'une adjonction (page 70)}

Une feuille isolée, numérotée 1 dans un cercle, qui reprend la construction
standard de la page 13 et démontre ce que cette page-là avait seulement
annoncé.

Soit une adjonction $f \dashv g$ avec $f : A \to B$, d'unité
$\eta : \mathrm{id}_{A} \to gf$ et de coünité
$\varepsilon : fg \to \mathrm{id}_{B}$, et soit $\varphi = fg$ la comonade
associée, de comultiplication $\delta = f * \eta * g$.\footnote{La page écrit
« $g$ adjoint à \emph{gauche} de $f$ », le mot étant récrit par-dessus un
autre. Les deux flèches qu'elle écrit à la ligne suivante disent l'inverse :
avec $\mathrm{id}_{A} \to gf$ et $fg \to \mathrm{id}_{B}$, c'est $f$ qui est
adjoint à gauche et $g$ adjoint à droite. La page 3, qui met en place la même
situation, écrit « droite » — également par-dessus un mot raturé. On a suivi
les formules.}

\medskip

\textbf{La construction.} Pour tout $X$ de $B$, les itérées $\varphi^{n+1}(X)$
forment un objet co-semi-simplicial $\Phi^{\bullet}(X)$, de terme
$\Phi^{n}(X) = \varphi^{n+1}(X)$, dont les cofaces
$\varphi^{n+1} \to \varphi^{n+2}$ sont les $n+1$ insertions de la
comultiplication,
\[
  \varphi^{k} * \delta * \varphi^{n-k} \qquad (0 \leqslant k \leqslant n) .
\]
\footnote{Le premier diagramme de la page porte, sur les deux flèches
$\varphi \rightrightarrows \varphi^{2}$, les étiquettes $\varphi * \beta$ et
$\beta * \varphi$, où $\beta$ est la coünité. Elles ne peuvent pas être justes
dans ce sens-là : la coünité fait descendre $\varphi^{2}$ sur $\varphi$. Ce
sont les insertions de $\lambda = \delta$ qui font monter, et c'est bien
$\lambda$ que le dernier diagramme de la page porte sur la même paire de
flèches. Le NB de la page confirme la lecture donnée ici : « les opérations de
dégénérescence n'utilisent que $\beta$ » — les codégénérescences, absentes d'un
objet co-semi-simplicial, sont celles qui emploient la coünité.}

\medskip

\textbf{Le seul fait qui compte.} La construction n'a d'intérêt que par sa
trivialité du bon côté. Pour tout $Y$ de $A$, l'objet $X = f(Y)$ porte une
coaugmentation
\[
  f(Y) \longrightarrow \Phi^{\bullet}\bigl( f(Y) \bigr) ,
\]
et cette coaugmentation est une \emph{équivalence d'homotopie}. L'homotopie
contractante est donnée par une seule flèche, l'unité poussée par $f$ :
\[
  f \xrightarrow{\;f * \eta\;} \varphi f = fgf .
\]
C'est ce qu'on appelle aujourd'hui une \emph{dégénérescence supplémentaire} :
un objet coaugmenté qui en possède une est contractile, sans autre
vérification. Il n'y a pas de coaugmentation pour un $X$ quelconque, et la page
le dit entre parenthèses : il faut $X \simeq f(Y)$.

\medskip

\textbf{Le retour à la page 3.} La page se termine en demandant quels sont, pour
$X$ quelconque, les morphismes d'objets co-semi-simpliciaux
$X \to \Phi^{\bullet}(X)$. La réponse est une flèche $u : X \to \varphi(X)$
rendant commutatif

\begin{tikzcd}[column sep=large]
X \arrow[r, "u"] &
\varphi(X) \arrow[r, bend left=18, "\varphi(u)"] \arrow[r, bend right=18, "\delta"'] &
\varphi^{2}(X)
\end{tikzcd}

c'est-à-dire exactement une structure de coalgèbre sur $X$, la condition étant
mot pour mot celle de la page 3. Le dossier revient à son point de départ : un
tapis cartésien, c'est une manière de coaugmenter la construction standard.

\pagerange{71}{75}
\section*{Fidélité : ce qu'est un foncteur d'oubli (pages 71 à 75)}

La chemise, page 71, porte un mot : « Fidélité ». Le texte tient en un titre,
qu'il faut lire comme une définition en attente : \emph{les foncteurs fidèles
transportables comme foncteurs d'oubli de structures}.

\medskip

\textbf{Le problème.} On dit couramment qu'un groupe est un ensemble muni d'une
structure, et que le foncteur qui envoie un groupe sur son ensemble sous-jacent
« oublie » cette structure. La question est de savoir quelle propriété d'un
foncteur $p : E \to B$ fait de lui un oubli de structure, sans supposer connue
à l'avance la notion de structure.

\medskip

\textbf{Transportabilité.} $p$ est dit \emph{transportable} lorsque le foncteur
induit $E \times_{B} B^{\mathrm{is}} \to B^{\mathrm{is}}$ est une fibration, où
$B^{\mathrm{is}}$ est le sous-groupoïde maximal de $B$. Cela revient à dire
qu'un isomorphisme $b \xrightarrow{\sim} b'$ de la base se relève : une
structure sur $b$ se transporte le long de lui. C'est le \emph{transport de
structure} au sens naïf, et tout foncteur est isomorphe à un foncteur
transportable.\footnote{Le premier tiers de la page 72 est récrit deux fois et
presque entièrement biffé ; ne subsistent que les énoncés. La page ajoute que
deux foncteurs transportables isomorphes $p, p' : E \to B$ sont reliés par une
auto-équivalence de $E$ isomorphe à l'identité, ce qui est la forme précise de
l'unicité du transport.}

\medskip

\textbf{Ce que les fibres disent.} Ce sont les deux énoncés qui font la liasse.
Soit $p$ une fibration.

\begin{itemize}
\item $p$ est \emph{fidèle} si et seulement si les catégories fibres $E_{b}$
sont des ensembles préordonnés — au plus une flèche entre deux objets.
\item $p$ est \emph{conservatif} si et seulement si les catégories fibres sont
des groupoïdes.
\item Donc $p$ est fidèle et conservatif si et seulement si les fibres sont
\emph{discrètes}, c'est-à-dire des ensembles nus.
\end{itemize}

\footnote{Le manuscrit énonce ces équivalences pour $p$ seulement
\emph{transportable}, et dit « catégories ordonnées » là où on a écrit
« préordonnées ». Un sens ne demande rien : $p$ fidèle force chaque fibre à
être un préordre, puisque toutes ses flèches ont même image. L'autre sens
demande de pouvoir factoriser une flèche quelconque par une flèche cartésienne,
donc que $p$ soit une fibration ; c'est l'hypothèse qu'on a explicitée. Quant à
l'antisymétrie, l'auteur y revient lui-même au bas de la page 74 : les fibres
sont en général préordonnées et non ordonnées, et il faut les réduire.}

\medskip

\textbf{Le théorème.} Pour un foncteur $p : E \to B$, les conditions suivantes
sont équivalentes :

\begin{enumerate}
\item[a)] $p$ est une fibration, fidèle et conservative ;
\item[b)] $p$ est une fibration à fibres discrètes ;
\item[c)] pour tout objet $X$ de $E$, le foncteur induit
$E_{/X} \to B_{/p(X)}$ est une équivalence ;
\item[d)] il existe un préfaisceau $F$ sur $B$ et une $B$-équivalence
$E \approx B_{/F}$.
\end{enumerate}

C'est le théorème des \emph{fibrations discrètes} : une catégorie au-dessus de
$B$ dont la projection est une fibration à fibres discrètes n'est rien d'autre
que la catégorie des éléments d'un préfaisceau.\footnote{$B_{/F}$ désigne la
catégorie des éléments de $F$ : objets les couples $(b, s)$ avec $s \in F(b)$,
flèches celles de $B$ compatibles. C'est le cas discret de la construction de
Grothendieck. La marge de la page 74 pose la question — « la deuxième condition
de c) signifie que $p$ est fibrant à fibres discrètes ? » — et la réponse est
oui.}

Et la conclusion, qui est la définition cherchée : \emph{les foncteurs d'oubli
de structure sont les foncteurs transportables et fidèles.} On n'a pas besoin
de la conservativité : les fibres sont alors préordonnées, non discrètes, et
c'est bien ce qu'on veut — sur un même objet de $B$, deux structures peuvent
être comparables sans être égales.

\medskip

\textbf{La reconstruction.} Partant d'un foncteur fidèle transportable, on
définit sur $B^{\mathrm{is}}$
\[
  b \longmapsto \text{classes d'isomorphie de } E_{b}
  = \text{« structures d'espèce } p \text{ sur } b \text{ »} ,
\]
et $E$ se reconstitue à équivalence près : un objet de $E$ est un objet de $B$
muni d'une structure d'espèce $p$. Mais le foncteur
$E^{\mathrm{is}} \to B^{\mathrm{is}}$ ne suffit pas à retrouver $E$ elle-même.
Pour cela il faut que $p$ soit une fibration, et définir sur $B$ la notion de
structure image inverse le long d'une flèche $b \to c$, ce qui donne un
foncteur
\[
  \Sigma : B^{\circ} \longrightarrow (\text{ens.\ ordonnés}) .
\]
La connaissance de $\Sigma$ détermine $E$ à équivalence près.\footnote{C'est la
version « à fibres préordonnées » du théorème précédent : les fibrations
fidèles sur $B$ correspondent aux préfaisceaux d'ensembles ordonnés, comme les
fibrations discrètes correspondent aux préfaisceaux d'ensembles. La page ajoute
qu'à isomorphisme près, et non seulement à équivalence près, il faut une
condition supplémentaire que la lecture ne rend pas.}

\pagerange{76}{80}
\section*{Questions sur les topos (pages 76 à 80)}

La chemise, page 76, porte : « Morphismes et comorphismes — Liste Problèmes ».
Ce qui suit est une liste de onze questions, adressée à quelqu'un — la première
ligne porte « (Pour M\textsuperscript{lle} \ldots{} ?) ». Quatre feuillets au
crayon, très surchargés ; ce sont les plus difficiles du lot, et plusieurs
questions ne se laissent lire qu'en partie. On donne celles qui se
tiennent.\footnote{Les marges portent, d'une encre différente, des réponses
venues plus tard : un « O.K.\ pour immersions ouvertes » en face de la question
1, un « non » en face de la question 7, un « Il n'est ni l'un ni l'autre » en
face de la question 2. C'est le seul endroit du dossier où l'on voie l'auteur
revenir sur ses propres questions.}

\begin{enumerate}
\item[1)] Définir et étudier les \emph{immersions} de topos. Les caractériser
par la pleine fidélité de $f_{*}$, c'est-à-dire par
$f^{*}f_{*} \simeq \mathrm{id}$ — condition qui fait de $f^{*}$ quelque chose
de voisin d'un passage au quotient. Dans quelle mesure une immersion détermine
l'ouvert ou le fermé qui lui donne naissance ? Si deux immersions ont même
« lieu », l'équivalence qui les identifie est-elle unique à isomorphisme unique
près ? Caractériser les immersions \emph{ouvertes} — par le fait que $f^{*}$
commute aux limites projectives quelconques, en plus de la pleine fidélité de
$f_{*}$ — et les immersions \emph{fermées} — par une condition de
surjectivité de $F \to f_{*}f^{*}F$.

\item[2)] Étudier les morphismes tels que $f^{*}$ admette un adjoint à gauche
$f_{!}$. Le cas $U \to U_{/S}$ est-il typique ? \emph{Non} : la marge répond
qu'il n'est ni l'un ni l'autre, et recommande de regarder plutôt la commutation
aux limites.\footnote{Un morphisme dont l'image inverse admet un adjoint à
gauche est ce qu'on appelle aujourd'hui un morphisme \emph{essentiel}. La
question rejoint celle du théorème de la page 6, dans le premier lot, où
l'essentialité est caractérisée par la commutation de $\varphi$ aux produits
infinis.}

\item[3)] Étudier le cas où $f^{*}$ est pleinement fidèle, et celui où $f^{*}$
commute aux limites projectives \emph{finies}.

\item[4)] Étudier le cas où $f_{*}$ est pleinement fidèle.

\item[5)] Notion de morphisme \emph{dominant} de topos, définie par
$f_{*}(\emptyset) = \emptyset$. Lien avec la conservativité de $f^{*}$. Que
dire du cas où $f_{*}$ est conservatif ?

\item[6)] Caractériser les foncteurs exacts $U \to V$ qui proviennent d'un
morphisme — c'est-à-dire donner une forme intrinsèque au théorème de Cartier et
Artin.\footnote{La page renvoie à Deligne. C'est la question de reconnaître,
parmi les foncteurs exacts entre topos, ceux qui sont des images inverses ; le
théorème de recollement des pages 54 à 68 en est le cas particulier de deux
morceaux.}

\item[7)] Les morphismes de $\mathcal{U}$-topos forment une
$\mathcal{U}$-catégorie : les $\mathrm{Hom}$ sont petits. L'ensemble des
classes d'isomorphie de tels morphismes est-il petit lui aussi ? La marge
répond : \emph{non}.

\item[8)] Théorie de l'adjonction

\begin{tikzcd}[column sep=huge]
(\mathrm{Top}) \arrow[r] &
(\mathrm{Cat}) \arrow[l, bend left=30, "C \mapsto \widehat{C}"]
\end{tikzcd}

entre topos et catégories, et étude des topos dont les $H^{i}$ sont finis.

\item[9)] Définir les \emph{catégories fibrées de topos}, et décrire la
situation dans certains cas.\footnote{C'est la question à laquelle le lot
suivant est consacré : les pages 82 à 95 traitent des catégories fibrées, des
sites fibrés et du topos total.}

\item[10)] Catégories fibrées et champs essentiellement représentables.

\item[11)] La chaîne

\begin{tikzcd}[column sep=huge]
(\mathrm{Ens.}) \arrow[r] &
(\mathrm{Esp.}) \arrow[r] &
(\mathrm{Top}) \arrow[ll, bend left=30]
\end{tikzcd}

et, en marge, l'indication : 2-adjoint du foncteur \emph{fibres}.
\end{enumerate}

\medskip

\textbf{Les chaînes d'adjoints.} La page 79, entièrement barrée, cherche un
exemple de foncteur qui soit à la fois adjoint à gauche et adjoint à droite du
même foncteur, et propose une petite catégorie explicite à deux objets $e$, $a$
et trois flèches non triviales $u : e \to a$, $v : a \to e$, $w = vu$, soumises
à $uv = \mathrm{id}_{e}$, $uw = u$, $wv = v$ — la composition n'étant pas
entièrement déterminée par ces relations, puisqu'il reste à décider ce que vaut
$w^{2}$.

La page 80 énonce alors le résultat positif. Étant donné une chaîne de
foncteurs adjoints successifs, on peut la prolonger indéfiniment vers le bas en
notant $f_{!}, f^{*}, f_{*}, f^{!}, \ldots$ les adjoints itérés :

\begin{tikzcd}[column sep=small, row sep=small, nodes={font=\scriptsize}]
f_{!} & f^{*} \arrow[d, no head] & f_{*} \arrow[d, no head] &
  f^{!} \arrow[d, no head] & f_{?} \arrow[d, no head] \\
& g_{!} & g^{*} \arrow[d, no head] & g_{*} \arrow[d, no head] &
  g^{!} \arrow[d, no head] \\
& & h_{!} & h^{*} & h_{*}
\end{tikzcd}

\textbf{Proposition.} Soit $\mathcal{U}$ un univers contenant un ensemble
infini. Pour tout $n$, il existe deux $\mathcal{U}$-topos $E$ et $F$ de la
forme $\widehat{C}$, $\widehat{C'}$ avec $C$, $C'$ \emph{petites}, et $n$
foncteurs $f_{1}, \ldots, f_{n}$ alternativement de $E$ dans $F$ et de $F$ dans
$E$, chacun adjoint à droite du précédent, dont aucun n'est pleinement
fidèle.\footnote{L'énoncé est très lacunaire sur la page — la clause « dont
aucun n'est pl.\ fid. » est une insertion interlinéaire, et le rôle du composé
$g = f_{1}\cdots f_{n}$ n'est pas lisible. On a donné ce que la page
soutient. Une telle suite est ce qu'on appelle aujourd'hui une \emph{chaîne
d'adjoints} de longueur $n$.}

\textbf{Question.} Peut-on trouver une suite infinie dans les deux directions ?
Il suffirait de trouver $f$ et $g$ tels que $g$ soit à la fois adjoint à gauche
et adjoint à droite de $f$.

La phrase s'arrête là, au milieu ; elle se termine page 81, dans le lot
suivant.
\pagerange{81}{81}
\section*{Questions sur les topos : la fin (page 81)}

Six lignes, qui achèvent la phrase interrompue au bas de la page 80. La
question était : peut-on prolonger indéfiniment, dans les deux sens, une chaîne
de foncteurs adjoints successifs ? Il suffirait pour cela de disposer de $f$ et
$g$ tels que $g$ soit à la fois adjoint à gauche et adjoint à droite de $f$.

L'ébauche de réponse passe par une description d'un type d'objets : des couples
formés d'un ensemble $A$ muni d'une involution $w$, et d'une rétraction du
quotient $A/w$ sur la partie fixe $A^{w} \subset A/w$. Le foncteur $f$ associe
alors, à un ensemble $E$, le triple $(E, \mathrm{id}_{E}, E, \mathrm{id}_{E})$.

C'est tout ce que la page soutient. Ni la construction ni la conclusion ne sont
lisibles, et on ne les a pas reconstituées.\footnote{Trois mots de ces six
lignes sont illisibles, dont le nom même de l'objet que $f$ produit. La
question elle-même reste donc ouverte sur la page. Une chaîne d'adjoints
infinie dans les deux sens est en fait impossible dans des situations très
générales — mais le dire ici serait prêter à l'auteur une réponse que sa page
n'a pas.}

\pagerange{82}{95}
\section*{Catégories fibrées (pages 82 à 95)}

Une chemise, page 82, porte le titre et un pointage d'archiviste de vingt-deux
pages ; les feuillets conservés n'en comptent que treize, dont deux versos
vierges. Trois textes suivent, chacun sous un titre de la main de l'auteur.

\medskip

\noindent\textbf{Rappel des définitions.} Soit $p : E \to B$ un foncteur. Une
flèche $\xi \to \eta$ de $E$ est \emph{cartésienne} si elle est universelle
parmi celles qui se projettent de la même façon ; $p$ est une \emph{fibration},
ou est \emph{fibrante}, si toute flèche $x \to y$ de $B$ et tout $\eta$ au-dessus
de $y$ admettent un relèvement cartésien. On note $E_{b}$ la fibre au-dessus de
$b$, et $f^{*} : E_{y} \to E_{x}$ le foncteur de transport attaché à
$f : x \to y$. Dualement, $p$ est \emph{cofibrante} si $p^{\circ}$ est
fibrante, et le transport va alors dans l'autre sens, $f_{*} : E_{x} \to
E_{y}$.

\pagerange{83}{84}
\subsection*{Sites fibrés et topos total}

Soit $E \to B$ une catégorie fibrée dont chaque fibre $E_{b}$ est munie d'une
topologie, les foncteurs de transport $f^{*}$ étant des morphismes de sites.
On munit alors $E$ de la topologie la moins fine rendant continus tous les
foncteurs $E_{b} \to E$ ; les familles couvrantes sont celles que décrit SGA 4
IV 3.2.\footnote{La page renvoie à « IV 3.2 » sans plus.}

Supposons $B$ et $E$ dans un univers $\mathcal{U}$. Le but est de calculer le
topos $\widetilde{E}$ des faisceaux sur $E$ à partir des topos $\widetilde{E}_{b}$
des fibres. Le premier pas consiste à plonger chaque fibre dans son topos,
c'est-à-dire à fabriquer une nouvelle catégorie fibrée $\overline{E}$ dont les
fibres sont les $\widetilde{E}_{b}$ : un morphisme de $F \in \overline{E}_{x}$
vers $G \in \overline{E}_{y}$ au-dessus de $f : x \to y$ est un morphisme
\[
  F \longrightarrow f^{*}(G) \quad \text{dans } \widetilde{E}_{x} ,
\]
où l'on note encore $f^{*}$ le prolongement du transport aux
topos.\footnote{Le prolongement existe parce que $f^{*}$ est un morphisme de
sites : il induit un morphisme de topos, dont l'image inverse prolonge $f^{*}$.
La page ne le dit pas.}

\medskip

\textbf{La formule.} Le topos total s'écrit alors comme une catégorie de
sections :
\[
  \widetilde{E}^{\circ} \;\simeq\;
  \underline{\Gamma}^{\,\text{vert-cart}}
  \bigl( B,\; X \mapsto (\overline{E})^{\circ}_{p(X)/X} \bigr) ,
\]
l'exposant « vert-cart » signifiant qu'on se borne aux sections cartésiennes
au-dessus des flèches verticales. Par associativité des sections, le second
membre se réécrit
\[
  \underline{\Gamma}\Bigl( B,\; x \mapsto
  \underline{\Gamma}^{\,\mathrm{cart}}\bigl( E_{x},\;
  X \mapsto (\overline{E}_{x})^{\circ}_{/X} \bigr) \Bigr) ,
\]
le foncteur intérieur n'étant autre que $\widetilde{E}^{\circ}_{x}$, d'où
enfin
\[
  \widetilde{E} \;\simeq\;
  \underline{\Gamma}\bigl( B, \overline{E}^{\Delta} \bigr)^{\circ} ,
\]
où $\overline{E}^{\Delta}$ se déduit de $\overline{E}$ par une opération sur
les fibres.\footnote{L'opération est illisible : le mot qui la nomme est
inséré au-dessus d'un mot biffé, et la transcription le donne comme
« groupoïdes » sans certitude. Une première formule, encadrée puis barrée sur
la page, proposait $\widetilde{E} \simeq \varprojlim_{X}(E_{p(X)})_{/X} =
\varprojlim_{x} E_{x}$ ; c'est elle que la version ci-dessus corrige, et la
correction porte précisément sur le fait qu'il faut des sections cartésiennes,
non une limite projective naïve.}

Le contenu, une fois les formules mises de côté, est celui-ci : \emph{un
faisceau sur la catégorie totale est une famille de faisceaux, un par fibre,
compatibles aux transports.} C'est la définition du topos total d'un site
fibré, et c'est la dernière apparition, dans le dossier, du motif « objet
global = famille locale + donnée de recollement ».

\pagerange{86}{88}
\subsection*{Sections et sections cartésiennes : les adjoints}

Le deuxième texte porte, encadré en haut de la page 86, le titre :
« Catégories $\underline{\Gamma}$ et $\varprojlim$ ».

\medskip

\noindent\textbf{Une notation à retenir.} Sur ces pages, et contrairement à
l'usage d'aujourd'hui, l'indice $*$ marque un adjoint à \emph{gauche} et
l'indice $+$ un adjoint à \emph{droite}. Les deux figurent côte à côte page 86,
et la page 87 conclut sur $F^{c}_{+}$ après avoir conclu sur $F^{c}_{*}$ ; la
distinction est donc voulue, et il faut la lire.\footnote{On a gardé les
indices de l'auteur, la page en dépendant. Le lecteur venu de la littérature
moderne, où $(-)_{*}$ est presque toujours un adjoint à droite, doit s'en
défier ici.}

\medskip

Pour une catégorie fibrée $E \to B$, on dispose de deux catégories : celle des
sections quelconques, notée $\underline{\Gamma}(E/B)$, et celle des sections
\emph{cartésiennes}, notée $\varprojlim E/B$ — ce sont les familles
compatibles aux transports, et c'est bien une limite projective, au sens
2-catégorique.\footnote{Aujourd'hui on écrit plutôt
$\lim_{B^{\circ}} E_{\bullet}$ ou $\mathrm{Cart}_{B}(B, E)$ ; la notation
$\varprojlim E/B$ du manuscrit dit la même chose.} L'inclusion
\[
  i : \varprojlim E/B \lhook\joinrel\longrightarrow \underline{\Gamma}(E/B)
\]
et un changement de base $B' \to B$ donnent le carré

\begin{tikzcd}
& \varprojlim E/B \arrow[r, hook, "i"] \arrow[d, "F^{*}_{c}"'] &
\underline{\Gamma}(E/B) \arrow[d, "F^{*}"] \\
X' \in B & \varprojlim E'/B' \arrow[r, hook, "i'"] &
\underline{\Gamma}(E'/B')
\end{tikzcd}

et toute la question est de savoir quand $F^{*}_{c}$, le changement de base
restreint aux sections cartésiennes, admet un adjoint.

\medskip

\textbf{Les deux calculs.} Ils tiennent en trois lignes chacun, et ils sont
symétriques. Du côté gauche,
\[
  \mathrm{Hom}\bigl( X', F^{*}_{c}(Y) \bigr)
  = \mathrm{Hom}\bigl( i'(X'), F^{*}(i(Y)) \bigr)
  = \mathrm{Hom}\bigl( F_{*} i'(X'), i(Y) \bigr)
  = \mathrm{Hom}\bigl( i_{*} F_{*} i'(X'), Y \bigr) ,
\]
ce qui identifie $F^{c}_{*} = i_{*} F_{*} i'$ ; du côté droit,
\[
  \mathrm{Hom}\bigl( F^{*}_{c}(Y), X' \bigr)
  = \mathrm{Hom}\bigl( F^{*}(i(Y)), i'(X') \bigr)
  = \mathrm{Hom}\bigl( i(Y), F_{+} i'(X') \bigr)
  = \mathrm{Hom}\bigl( Y, i_{+} F_{+} i'(X') \bigr) ,
\]
ce qui identifie $F^{c}_{+} = i_{+} F_{+} i'$. D'où les deux énoncés :

\begin{enumerate}
\item[\textbf{①}] Si $i$ admet un adjoint à gauche $i_{*}$ et $F^{*}$ un
adjoint à gauche $F_{*}$ — ce qui a lieu par exemple si $E$ est assez cofibrée
sur $B$ et si les « petites » limites projectives existent dans les fibres —
alors $F^{*}_{c}$ admet un adjoint à gauche $F^{c}_{*}$.
\item[\textbf{②}] Si $i$ admet un adjoint à droite $i_{+}$ et $F^{*}$ un
adjoint à droite $F_{+}$ — par exemple si les petites limites projectives
existent dans les fibres de $E/B$ — alors $F^{*}_{c}$ admet un adjoint à droite
$F^{c}_{+}$.
\end{enumerate}

\medskip

\textbf{③ Proposition.} Supposons que pour tout $b$ de $B$ la limite inductive
\[
  \varinjlim_{c \in \mathrm{Ob}\,{}^{b}\!/B} f^{*}\bigl( X(c) \bigr)
\]
existe dans $E_{b}$, et qu'elle commute à tous les changements de base
$u^{*}$. Alors l'adjoint à gauche du foncteur d'inclusion existe, et il est
donné par cette formule même :
\[
  \dot{F}^{\varepsilon}_{*}(X)(b)
  = \varinjlim_{x \in \mathrm{Ob}\,{}^{b}\!/B} f^{*}\bigl( X(x) \bigr) .
\]
Autrement dit : \emph{pour rendre cartésienne une section quelconque, on
recolle fibre à fibre, par une limite inductive sur les objets au-dessus de
$b$.}\footnote{Le nom du foncteur porte sur la page un point et un exposant
qu'on n'a pas su lire ; on a gardé la graphie de la transcription. La page
écrit par ailleurs le but de l'inclusion $\underline{\Gamma}(B/E)$, alors que
partout ailleurs elle écrit $\underline{\Gamma}(E/B)$. Le dernier tiers de la
page 87, une dizaine de lignes et une note marginale encadrée, est barré d'un
long trait diagonal et reste illisible.}

\medskip

\textbf{Deux corollaires} (page 88). Le premier : si de plus les $b_{*}$
existent — par exemple si $p$ est aussi cofibrant et si les petites sommes
existent dans les fibres —, alors les « foncteurs fibres »
$b^{*}_{c} : \varprojlim E/B \to E_{b}$ admettent chacun un adjoint à gauche
$b^{c}_{*}$.

Le second concerne les générateurs, et c'est celui qui sert. Si
$\underline{\Gamma}(E/B)$ a une famille de générateurs stricts
$(X_{\lambda})$, alors les $i(X_{\lambda})$ en forment une pour $\varprojlim
E/B$. En particulier, si $p$ est cofibrant et si chaque fibre $E_{b}$ a des
petites sommes et une petite famille de générateurs stricts, alors $\varprojlim
E/B$ en a une aussi, obtenue en réunissant celles des fibres :
\[
  \bigcup_{b \in \mathrm{Ob}\,B} b^{c}_{*}\bigl( \mathcal{G}_{b} \bigr) ,
\]
où $\mathcal{G}_{b}$ est un système de générateurs stricts de
$E_{b}$.\footnote{La page écrit « en prenant pour la \ldots{} $E_{b}$ des gén.\
stricts de $E_{b}$ » ; le membre de phrase est incomplet, mais la formule
finale ne laisse pas de doute sur ce qui est réuni.}

\pagerange{89}{92}
\subsection*{Générateurs, cogénérateurs, et quand $\varprojlim E/B$ est un topos}

Quatre pages au crayon, lourdement surchargées, dont la transcription est plus
lacunaire que pleine. On donne ce qui est établi, et on s'arrête là.

\medskip

Le fil est celui-ci. Pour obtenir un système de générateurs de
$\varprojlim E/B$, il faut passer par l'adjoint à gauche $i_{*}$ de
l'inclusion, et non par l'inclusion elle-même : si $(X_{\lambda})$ engendre
$\underline{\Gamma}(E/B)$, la bijection
\[
  \mathrm{Hom}_{\underline{\Gamma}}\bigl( X_{\lambda}, i(Y) \bigr)
  \simeq \mathrm{Hom}\bigl( i_{*}(X_{\lambda}), Y \bigr)
\]
montre que ce sont les $i_{*}(X_{\lambda})$ qui engendrent en bas. Pour des
\emph{co}générateurs, c'est-à-dire des générateurs de la catégorie opposée,
c'est l'adjoint de l'autre côté qu'il faut, et la page 90 en tire une condition
suffisante : que cet adjoint existe et qu'il y ait un petit système de
générateurs dans $\underline{\Gamma}(E^{\Delta}/B)$ — ce qui a lieu, par
exemple, si les petites limites projectives existent dans les fibres, si les
foncteurs de transport y commutent, et si les fibres ont de petits ensembles de
cogénérateurs.

\medskip

Les pages 91 et 92 posent la question finale : quand $\varprojlim E/B$ est-elle
un $\mathcal{U}$-topos ? Le principe est un transfert : si un type de limite
est représentable dans chaque fibre et si les transports y commutent, alors ce
type de limite est représentable dans $\varprojlim E/B$ et les foncteurs fibres
y commutent. On en déduit que si les fibres sont préadditives, additives ou
abéliennes — et si les transports respectent ces structures —, il en va de même
de $\varprojlim E/B$. Reste, pour conclure au topos, deux conditions
supplémentaires : que la catégorie soit une $\mathcal{U}$-catégorie, ce qui
tient à la petitesse de $B$, et qu'elle possède une petite famille de
générateurs, ce que le corollaire de la page 88 fournit.\footnote{Les pages 91
et 92 poursuivent en discutant ce que coûte l'hypothèse de cofibration et le
rôle des adjoints $f_{!}$, mais elles sont illisibles par larges plages : leurs
interlignes sont plus denses que leurs lignes. La page 92 s'achève sur un aveu
— le calcul en général est « un peu fastidieux ». On n'a rien reconstitué de ce
qui manque.}

\pagerange{94}{95}
\subsection*{Dualité : fibré et cofibré}

Le dernier texte porte, encadré en haut de la page 94 : « Dualité pour
catégories fibrées et cofibrées ». Il tient en deux pages et il est clair d'un
bout à l'autre.

\medskip

\textbf{Les deux notions.} Soit $p : E \to B$.

\begin{enumerate}
\item[1)] $p$ est \emph{fibrante} si, pour toute flèche $f : x \to y$ de $B$,
un foncteur de transport $f^{*} : E_{y} \to E_{x}$ est défini, avec
transitivité essentielle, par l'identité
\[
  \mathrm{Hom}_{f}(X, Y) \simeq
  \mathrm{Hom}_{\mathrm{id}_{x}}\bigl( X, f^{*}(Y) \bigr) .
\]
De telles données reviennent à un pseudo-foncteur $F : B^{\circ} \to
(\mathrm{Cat})$.
\item[2)] $p$ est \emph{cofibrante}, c'est-à-dire $p^{\circ}$ fibrante, si pour
toute flèche $f : x \to y$ un foncteur $f_{*} : E_{x} \to E_{y}$ est défini,
avec transitivité essentielle, par
\[
  \mathrm{Hom}_{f}(X, Y) \simeq
  \mathrm{Hom}_{\mathrm{id}_{y}}\bigl( f_{*}(X), Y \bigr) .
\]
\footnote{Le second membre est d'une lecture douteuse sur la page ; la
transcription le donne comme $\mathrm{Hom}\,\varphi_{p}(X', Y)$. C'est la
symétrie avec 1) qui impose la forme écrite ici, et c'est elle que la suite du
texte emploie.}
De telles données reviennent à un pseudo-foncteur $F : B \to (\mathrm{Cat})$.
\end{enumerate}

\medskip

\textbf{Le lien.} C'est l'énoncé qui fait la page, et il est purement
adjonctif.

\begin{itemize}
\item Si $p$ est fibrante : pour une flèche $f$ donnée, $f_{*}$ existe si et
seulement si $f^{*}$ admet un adjoint à gauche. Donc $p$ est cofibrante si et
seulement si \emph{tous} les $f^{*}$ admettent un adjoint à gauche.
\item Si $p$ est cofibrante : pour une flèche $f$ donnée, $f^{*}$ existe si et
seulement si $f_{*}$ admet un adjoint à droite. Donc $p$ est fibrante si et
seulement si tous les $f_{*}$ admettent un adjoint à droite.
\end{itemize}

Une catégorie à la fois fibrée et cofibrée est donc exactement une famille de
catégories dont tous les transports vont par paires adjointes. C'est le même
énoncé que celui qui ouvre le dossier, à quatre-vingt-dix pages de distance :
tout se lit sur les adjonctions.

\medskip

\textbf{L'avertissement final.} La \emph{catégorie opposée relative} d'une
catégorie fibrée $p : E \to B$ — celle qu'on obtient en retournant chaque
fibre — n'est définie qu'à équivalence près. C'est une autre catégorie fibrée,
associée au pseudo-foncteur
\[
  F^{\mathrm{opp}} : B \longrightarrow (\mathrm{Cat}) , \qquad
  F^{\mathrm{opp}}(a) = F(a)^{\circ} , \qquad
  F^{\mathrm{opp}}(f) = F(f)^{\circ} .
\]
Il n'est pas évident, note l'auteur, qu'on puisse en donner une description
intrinsèque, qui évite le recours au pseudo-foncteur. Et pour que
$(E/B)^{\mathrm{opp}}$ soit à son tour cofibrée, il faut et il suffit que les
$f^{*}$ admettent des adjoints à \emph{droite} — condition qui n'est autre que
celle pour que $E/B$ soit cofibrée.

La dernière phrase du dossier est celle-ci, et elle est soulignée sur la page :
\emph{à retenir que ce pseudo-foncteur définit deux catégories fibrées bien
distinctes sur $B$.}\footnote{La variance est ce qui les distingue : un même
$F$ peut être lu comme pseudo-foncteur sur $B^{\circ}$, donnant une fibration,
ou sur $B$, donnant une cofibration, et les deux catégories totales n'ont
aucune raison de coïncider. C'est la confusion contre laquelle l'auteur met en
garde, et c'est la note sur laquelle les quatre-vingt-quinze pages
s'arrêtent.}

\end{document}
